\documentclass[11pt,a4paper]{article}

% ============================================================
% Uniform MTDF preamble -- identical across all papers
% ============================================================
\usepackage[utf8]{inputenc}
\usepackage[T1]{fontenc}
\usepackage{lmodern}
\usepackage[english]{babel}
\usepackage[a4paper,margin=2.2cm]{geometry}
\usepackage{amsmath,amssymb,amsthm}
\usepackage{graphicx}
\usepackage{float}
\usepackage{booktabs}
\usepackage{array}
\usepackage{longtable}
\usepackage{tabularx}
\usepackage{xltabular}
\usepackage{multirow}
\usepackage{caption}
\usepackage[dvipsnames,table]{xcolor}
\usepackage{mdframed}
\usepackage{parskip}
\usepackage{ragged2e}
\usepackage[hidelinks,breaklinks=true]{hyperref}
\usepackage{microtype}
\usepackage{url}
% Allow \path{} to break more aggressively inside narrow table columns:
\makeatletter
\g@addto@macro\UrlBreaks{\do\_\do\-\do\.\do\/\do\:}
\makeatother

% Colours matching the HTML theme
\definecolor{accent}{HTML}{1A4B8A}
\definecolor{callout-bg}{HTML}{FBFBFB}
\definecolor{tablehead}{HTML}{F0F0F0}
\definecolor{retire}{HTML}{8A1A1A}

% Prevent any table from overflowing
\setlength{\tabcolsep}{4pt}
\renewcommand{\arraystretch}{1.15}

% Caption style (compact, italic, small like the HTML)
\captionsetup{font=small,labelfont=bf,labelsep=period,justification=centerfirst}

% Hyperref metadata placeholders (filled per-paper below)
\hypersetup{
  pdftitle={MTDF Independent GPU Validation: Phases 1–6},
  pdfauthor={Ingo Mesche},
  pdfsubject={MTDF - Mesche Time-Dilation Field Theory},
  colorlinks=false
}

% Prevent overfull hbox from long URLs/DOIs in bibliography
\sloppy
\emergencystretch=3em

% Tolerant line breaking so long identifiers (Python filenames etc.) don't
% produce large overfull hboxes in narrow table columns. These settings keep
% all content on-page while slightly relaxing right-margin strictness.
\tolerance=1200
\hbadness=2000
\hyphenpenalty=150
\exhyphenpenalty=50

% ============================================================
\title{MTDF Independent GPU Validation: Phases 1--6}
\author{Ingo Mesche\\ \small Independent Researcher, Malta \\ \small \texttt{imesche@outlook.com}}
\date{V75 / July 2026}

\begin{document}
\maketitle
\begin{center}\itshape Blind reproduction, CMB consistency, environment signal reproduction, sensitivity forecasts, full Planck hard-falsifier test, and cross-probe discriminatorsHardware: NVIDIA RTX 4080 SUPER (16 GB VRAM, CUDA 12.9)\end{center}



\begin{mdframed}[linecolor=accent,linewidth=0.5pt,backgroundcolor=callout-bg,leftline=true,rightline=true,topline=true,bottomline=true,innerleftmargin=8pt,innerrightmargin=8pt,innertopmargin=6pt,innerbottommargin=6pt]
\textbf{Abstract.}
This document reports the results of a fully independent, GPU-accelerated validation of Mesche's Tensor Dynamics Framework (MTDF). Six phases were executed over several weeks using code written from scratch, with no shared libraries or routines from the original MTDF analysis pipeline. Phase~1 reproduces all V18 dashboard $\chi$$^{2}$ values to four decimal places. Phase~2 tests MTDF against the Planck 2018 CMB power spectrum via MCMC with a perturbative correction layer, finding $\Delta$$\chi$$^{2}$~<~1 and an unconstrained early-field-energy amplitude. Phase~3 independently reproduces the SN~$\times$~void environment signal reported in MTDF\_\allowbreak{}02, with bootstrap confidence intervals excluding zero. Phase~4 forecasts that 5$\sigma$ detection of the environment signal requires approximately 3,400 low-redshift supernovae, placing definitive confirmation within the reach of LSST/\allowbreak{}Rubin. Phase~5 executed the Priority~1 hard falsifier, the full Planck 2018 plik TTTEEE + lowl + lensing likelihood via the \texttt{class\_\allowbreak{}mtdf} Boltzmann solver; its v1 numerical results were later retired as a documented gauge artifact, and the sealed v2 chains (corrected covariant $\mu$(a), newtonian gauge) now govern: standard parameters essentially unshifted, a broad coupling posterior k\textsubscript{f}~=~0.685~$\pm$~0.500 (consistency, not support), an enhanced-growth posture with $\sigma$\textsubscript{8}~=~0.8265~$\pm$~0.0062 about 1.9$\sigma$ (posterior units) above the $\Lambda$CDM baseline, and a completed comparison against the matched $\Lambda$CDM control chain: $\Delta$$\chi$$^{2}$ (bounded best fits) = $-$1.29. This is a maximum-likelihood comparison, not a Bayes factor or Bayesian evidence calculation. Phase~6 extends validation to cross-probe discriminators: a redshift transition scan (3.6$\sigma$ at z~<~0.03), weak lensing $\times$ environment (KiDS-1000; systematics-clean null with explicit upper limits), derived parameter consistency (v1-posterior-based; superseded, see 6b), and CMB lensing $\times$ voids (Planck~PR4 stacking on DESI~BGS and BOSS~DR12 void catalogues; pipeline validated against published detections).
\end{mdframed}

\tableofcontents

\begin{mdframed}[linecolor=accent,linewidth=0.5pt,backgroundcolor=callout-bg,leftline=true,rightline=true,topline=true,bottomline=true,innerleftmargin=8pt,innerrightmargin=8pt,innertopmargin=6pt,innerbottommargin=6pt]
\textbf{Claim status.}
This paper distinguishes between:
(i) calibration anchors,
(ii) validation targets,
(iii) diagnostic consistency checks, and
(iv) speculative extensions.
Only results explicitly labelled as validation targets are used as evidence for the core MTDF claim. Diagnostic and speculative sections are included to define future tests, not to claim established confirmation. Throughout, ``validated'' denotes independent reproduction within this programme (new code, same data); it does not denote external replication on independent data, which no MTDF result yet has.
\end{mdframed}

\section{Scope and methodology}
\label{sec:sec1}
The purpose of this validation exercise is to subject MTDF to a battery of independent tests that are entirely separate from the original analysis pipeline. All code was written from scratch for this programme. No libraries, routines, or intermediate results from the original MTDF codebase were used for Phases~1, 3, 4, or 6. Phase~2 used the CosmoPower $\Lambda$CDM emulator with a bespoke MTDF correction layer. Phase~5 used the existing \texttt{class\_\allowbreak{}mtdf} modified Boltzmann solver through cobaya with the full Planck likelihood stack.

The six phases test qualitatively different aspects of the framework:

\begin{itemize}
  \item \textbf{Phase 1} tests internal consistency: can the published $\chi$$^{2}$ values be reproduced from raw data using independent code?
  \item \textbf{Phase 2} tests CMB compatibility (perturbative): is MTDF allowed by a compressed CMB likelihood?
  \item \textbf{Phase 3} tests the primary empirical prediction: is the SN~$\times$~void environment signal real and reproducible?
  \item \textbf{Phase 4} tests future viability: what survey specifications are needed for definitive detection?
  \item \textbf{Phase 5} tests the hard falsifier: does MTDF survive the full Planck likelihood with modified Boltzmann equations?
  \item \textbf{Phase 6} tests cross-probe discriminators: does the transition signature survive in weak lensing, CMB lensing, and derived parameters?
\end{itemize}

\begin{mdframed}[linecolor=accent,linewidth=0.5pt,backgroundcolor=callout-bg,leftline=true,rightline=true,topline=true,bottomline=true,innerleftmargin=8pt,innerrightmargin=8pt,innertopmargin=6pt,innerbottommargin=6pt]
\textbf{Independence guarantee.}
Phases~1, 3, and 6 use fully independent implementations. Phase~2 uses a perturbative correction layer on top of a public $\Lambda$CDM emulator (CosmoPower), parameterised by a continuous coupling \( k_f \) where \( k_f = 0 \) is $\Lambda$CDM and \( k_f = 1 \) is full MTDF. Phase~4 builds on Phase~3 infrastructure with additional Monte Carlo machinery. Phase~5 uses the existing \texttt{class\_\allowbreak{}mtdf} Boltzmann solver with cobaya and the full Planck plik likelihood. This is the only phase that uses pre-existing MTDF code, which is appropriate since the test is specifically to validate that code against Planck. Phase~6 uses public data products (KiDS-1000 shear catalogues, Planck~PR4 lensing maps, DESI/\allowbreak{}BOSS void catalogues) with all stacking and analysis code written from scratch. All GPU computations were cross-checked against CPU results.
\end{mdframed}

\subsection{Hardware and software}
All computations were performed on a desktop workstation equipped with an NVIDIA RTX 4080 SUPER GPU (16 GB VRAM, CUDA 12.9, compute capability 8.9). The software environment consists of Python 3.12 with cobaya 3.6, emcee, getdist, CuPy, JAX, TensorFlow, and CosmoPower. All datasets were pre-downloaded: Pantheon+ with full covariance matrix, DESI Y1 BAO, cosmic chronometers, DR16 f$\sigma$\textsubscript{8} growth measurements, and DESI void catalogues (VoidFinder, REVOLVER, VIDE).

\section{Phase 1: Independent $\chi$$^{2}$ reproduction}
\label{sec:sec2}
PASSED: All V18 dashboard values reproduced to $\delta$ = 0.000049A fully independent implementation reproduced all V18 dashboard $\chi$$^{2}$ values. The code reads raw data files, constructs likelihood functions from scratch, and evaluates MTDF predictions at the V18 parameter point. No dashboard code was consulted or reused.

\subsection{Strict total}
\begin{table}[H]
\centering
\footnotesize
\begin{tabularx}{\linewidth}{l>{\RaggedRight\arraybackslash}X>{\RaggedRight\arraybackslash}X>{\RaggedRight\arraybackslash}X}
\toprule
\textbf{Metric} & \textbf{Target (V18)} & \textbf{Reproduced} & \textbf{$\delta$} \tabularnewline
$\chi$$^{2}$/$\nu$ & 1.1683 & 1.1683 & 0.000049 \tabularnewline
Degrees of freedom $\nu$ & 1745 & 1745 & -- \tabularnewline
\bottomrule
\end{tabularx}
\end{table}

Note (July 2026): the targets above are the V18-era values that this reproduction run verified; the current sealed totals (V19 workbook, V75 dashboard) are given in MTDF\_\allowbreak{}01, Section~7.

\subsection{Vector pillar breakdown}
\begin{table}[H]
\centering
\footnotesize
\begin{tabularx}{\linewidth}{l>{\RaggedRight\arraybackslash}X>{\RaggedRight\arraybackslash}X>{\RaggedRight\arraybackslash}X}
\toprule
\textbf{Pillar} & \textbf{$\chi$$^{2}$} & \textbf{DOF} & \textbf{$\chi$$^{2}$/$\nu$} \tabularnewline
Pantheon+ SNe Ia & 2012.72 & 1700 & 1.184 \tabularnewline
DESI Y1 BAO & 15.39 & 12 & 1.282 \tabularnewline
Cosmic chronometers H(z) & 6.98 & 15 & 0.466 \tabularnewline
DR16 f$\sigma$\textsubscript{8} growth & 2.00 & 3 & 0.667 \tabularnewline
CMB distance prior (diagnostic only) & 595.33 & 3 & 198.44 \tabularnewline
\multicolumn{4}{l}{\textbf{Note on the diagnostic value.} $\chi$$^{2}$/$\nu$ = 198 against the Planck compressed distance prior is expected, not anomalous: the prior fixes the shape parameters under a $\Lambda$CDM template, and any theory modifying the early-universe geometry (here via f\textsubscript{kick}) or late-time growth (here via $\mu$(a)) is biased away from the prior centre even when the uncompressed C\textsubscript{$\ell$} spectra agree well. The honest test is the full Planck likelihood (Section 6, sealed v2 chains): standard parameters essentially unshifted, with the theory value k\textsubscript{f} = 1 well inside the posterior; the completed comparison against the matched $\Lambda$CDM control is reported in Section 6.3. The compressed prior is reported here strictly as a diagnostic and excluded from strict totals.} \tabularnewline
\bottomrule
\end{tabularx}
\end{table}

Plus 15 scalar pillars contributing $\chi$$^{2}$ = 1.67, matching exactly.

\begin{mdframed}[linecolor=accent,linewidth=0.5pt,backgroundcolor=callout-bg,leftline=true,rightline=true,topline=true,bottomline=true,innerleftmargin=8pt,innerrightmargin=8pt,innertopmargin=6pt,innerbottommargin=6pt]
\textbf{Significance.} Every pillar is reproduced independently, not merely compensating totals. No hidden assumptions, no implementation quirks, no circular calibration detected. The V18 strict protocol is validated by independent audit.
\end{mdframed}

\section{Phase 2: Planck CMB consistency}
\label{sec:sec3}
PASSED: $\Delta$$\chi$$^{2}$ < 1; \, k\textsubscript{f} = 1 within 95\% CI in all runsPhase 2 tests whether MTDF is compatible with the Planck 2018 CMB power spectrum. The approach uses the CosmoPower neural network emulator for $\Lambda$CDM C\textsubscript{$\ell$} spectra, augmented by an analytic MTDF correction layer that modifies the sound horizon, angular diameter distance, and ISW contribution as functions of a continuous coupling parameter \( k_f \). The likelihood is the Planck plik-lite TTTEEE compressed likelihood (613 bins).

\subsection{Emulator validation}
\begin{table}[H]
\centering
\footnotesize
\begin{tabularx}{\linewidth}{l>{\RaggedRight\arraybackslash}X>{\RaggedRight\arraybackslash}X>{\RaggedRight\arraybackslash}X}
\toprule
\textbf{Spectrum} & \textbf{Prediction time} & \textbf{Mean error vs CAMB} & \textbf{Max error} \tabularnewline
TT & 7.5 ms & 0.09\% & 0.20\% \tabularnewline
TE & 4.1 ms & 0.46\% & 15.4\%* \tabularnewline
EE & 7.6 ms & 0.11\% & 0.29\% \tabularnewline
\bottomrule
\end{tabularx}
\end{table}

*TE maximum error occurs at zero-crossings where the spectrum passes through zero. This is expected and numerically harmless.

\subsection{Fixed-parameter comparison}
\begin{table}[H]
\centering
\footnotesize
\begin{tabularx}{\linewidth}{l>{\RaggedRight\arraybackslash}X>{\RaggedRight\arraybackslash}X>{\RaggedRight\arraybackslash}X}
\toprule
\textbf{Cosmology} & \textbf{$\chi$$^{2}$} & \textbf{Bins} & \textbf{$\chi$$^{2}$/bin} \tabularnewline
$\Lambda$CDM (CAMB reference) & 584.45 & 613 & 0.953 \tabularnewline
MTDF at k\textsubscript{f} = 1 (corrected) & 585.26 & 613 & 0.955 \tabularnewline
$\Delta$$\chi$$^{2}$ & +0.81 &  &  \tabularnewline
\bottomrule
\end{tabularx}
\end{table}

$\Lambda$CDM $\chi$$^{2}$ = 584.45 matches the Planck 2018 published value ($\sim$583). Pipeline correctly wired.

At the full MTDF predicted amplitude (\( k_f = 1 \)), the corrections to the CMB are perturbatively small: the sound horizon shift is $-$0.071\%, the angular scale shift is +0.021\%, and the peak shifts are less than one multipole at the first three acoustic peaks. The $\Delta$$\chi$$^{2}$ of +0.81 is far below any statistical significance threshold.

\subsection{Normalisation correction}
During development, a normalisation inconsistency was identified and corrected. Three different \( k_f \) conventions existed in the codebase (see Appendix). After correction, the physical mapping is:

\[
f_{\rm kick} = \frac{\lambda_{\rm MTDF}}{24} = \frac{(1 - \beta_{\rm eos})^2}{(1+\alpha)\times 24} = 0.003303
\]

At \( k_f = 1 \): $\theta$\textsubscript{s} ratio = 1.000212 (shift of +0.021\%), ISW boost = 0.25\%. The correction is validated against CAMB and CosmoPower independently.

\subsection{Production MCMC results}
Three production MCMC runs were completed (5000 steps, 32 walkers each), using the emcee affine-invariant sampler with GPU-accelerated likelihood evaluation.

\subsubsection{Model comparison}
\begin{table}[H]
\centering
\footnotesize
\begin{tabularx}{\linewidth}{l>{\RaggedRight\arraybackslash}X>{\RaggedRight\arraybackslash}X>{\RaggedRight\arraybackslash}X}
\toprule
\textbf{Metric} & \textbf{Planck-only} & \textbf{Planck + BAO + SN} & \textbf{Planck + BAO + SN (k\textsubscript{f}\,<\,2)} \tabularnewline
N\textsubscript{data} & 613 & 2,326 & 2,326 \tabularnewline
$\Lambda$CDM $\chi$$^{2}$ & 583.10 & 2,607.80 & 2,607.80 \tabularnewline
MTDF best-fit $\chi$$^{2}$ & 582.24 & 2,607.10 & 2,606.89 \tabularnewline
$\Delta$$\chi$$^{2}$ & $-$0.85 & $-$0.70 & $-$0.91 \tabularnewline
$\Delta$AIC & +1.15 & +1.30 & +1.09 \tabularnewline
\bottomrule
\end{tabularx}
\end{table}

$\Delta$$\chi$$^{2}$ negative = MTDF marginally better in raw fit. $\Delta$AIC positive = $\Lambda$CDM mildly preferred by Occam penalty for the extra k\textsubscript{f} parameter. Result: completely inconclusive; neither model is favoured.

\subsubsection{k\textsubscript{f} posterior}
\begin{table}[H]
\centering
\footnotesize
\begin{tabularx}{\linewidth}{l>{\RaggedRight\arraybackslash}X>{\RaggedRight\arraybackslash}X>{\RaggedRight\arraybackslash}X}
\toprule
\textbf{} & \textbf{Planck-only [0,\,5]} & \textbf{Combined [0,\,5]} & \textbf{Combined [0,\,2]} \tabularnewline
Best-fit & 4.63 & 0.085 & 0.086 \tabularnewline
Median & 2.63 & 1.93 & \textbf{0.93} \tabularnewline
68\% CI & {}[0.89, 4.28] & {}[0.59, 3.72] & {}[0.28, 1.65] \tabularnewline
95\% UL & 4.77 & 4.53 & 1.89 \tabularnewline
k\textsubscript{f}\,=\,1 in 95\%? & Yes & Yes & Yes \tabularnewline
k\textsubscript{f}\,=\,0 in 68\%? & Yes & Yes & Yes \tabularnewline
\bottomrule
\end{tabularx}
\end{table}

\subsubsection{H\textsubscript{0} posterior}
\begin{table}[H]
\centering
\footnotesize
\begin{tabularx}{\linewidth}{l>{\RaggedRight\arraybackslash}X>{\RaggedRight\arraybackslash}X>{\RaggedRight\arraybackslash}X}
\toprule
\textbf{} & \textbf{Planck-only} & \textbf{Combined} & \textbf{Combined (narrow)} \tabularnewline
H\textsubscript{0} (km/\allowbreak{}s/\allowbreak{}Mpc) & 66.89 $\pm$ 0.62 & 67.62 $\pm$ 0.40 & 67.64 $\pm$ 0.40 \tabularnewline
\bottomrule
\end{tabularx}
\end{table}

H\textsubscript{0} is rock-stable across all runs and independent of k\textsubscript{f} prior choice.

\subsubsection{Per-probe breakdown (combined narrow best-fit)}
\begin{table}[H]
\centering
\footnotesize
\begin{tabularx}{\linewidth}{l>{\RaggedRight\arraybackslash}X>{\RaggedRight\arraybackslash}X>{\RaggedRight\arraybackslash}X}
\toprule
\textbf{Probe} & \textbf{$\chi$$^{2}$} & \textbf{N\textsubscript{data}} & \textbf{$\chi$$^{2}$/N} \tabularnewline
Planck plik-lite TTTEEE & 582.74 & 613 & 0.951 \tabularnewline
DESI Y1 BAO & 20.27 & 12 & 1.689 \tabularnewline
Pantheon+ SNe Ia & 2,003.44 & 1,701 & 1.178 \tabularnewline
\textbf{Total} & \textbf{2,606.89} & \textbf{2,326} & \textbf{1.121} \tabularnewline
\bottomrule
\end{tabularx}
\end{table}

\begin{mdframed}[linecolor=accent,linewidth=0.5pt,backgroundcolor=callout-bg,leftline=true,rightline=true,topline=true,bottomline=true,innerleftmargin=8pt,innerrightmargin=8pt,innertopmargin=6pt,innerbottommargin=6pt]
\textbf{Phase 2 conclusion.}
MTDF is fully compatible with Planck 2018 TTTEEE + DESI BAO + Pantheon+ SNe. The predicted stress-field correction (\~{}0.13\% on H(z)) is below current data sensitivity. Full MTDF (\( k_f = 1 \)) and $\Lambda$CDM (\( k_f = 0 \)) are both allowed. Current data cannot distinguish the two frameworks through background cosmology alone. With the physics-motivated narrow prior [0,\,2], the posterior median is \( k_f = 0.93 \), essentially the theoretical prediction. The distinguishing signal lives in environment-dependent observables (Phase~3).
\end{mdframed}

\section{Phase 3: SN $\times$ void environment signal}
\label{sec:sec4}
PASSED: Exact reproduction; bootstrap 95\% CI excludes zeroPhase 3 independently reproduces the SN~Ia~$\times$~void environment analysis reported in MTDF\_\allowbreak{}02 (Companion Part~I). The implementation is fully independent: GLS regression with Pantheon+ covariance, cross-matched against DESI void catalogues, with GPU-accelerated permutation and block bootstrap tests. Runtime: 8.2 minutes (10,000 permutations + 5,000 bootstraps).

\subsection{Environment coefficient $\gamma$\textsubscript{env}}
\begin{table}[H]
\centering
\footnotesize
\begin{tabularx}{\linewidth}{l>{\RaggedRight\arraybackslash}X>{\RaggedRight\arraybackslash}X>{\RaggedRight\arraybackslash}X>{\RaggedRight\arraybackslash}X>{\RaggedRight\arraybackslash}X}
\toprule
\textbf{Catalogue} & \textbf{$\gamma$\textsubscript{env} (mag)} & \textbf{$\Delta$$\chi$$^{2}$} & \textbf{p-value} & \textbf{Target $\Delta$$\chi$$^{2}$} & \textbf{Target p} \tabularnewline
VoidFinder & +0.0030 $\pm$ 0.0018 & 2.92 & 0.088 & 2.91 & 0.088 \tabularnewline
REVOLVER & +0.0047 $\pm$ 0.0023 & 4.25 & 0.039 & 4.25 & 0.039 \tabularnewline
VIDE & +0.0046 $\pm$ 0.0026 & 3.15 & 0.076 & 3.15 & 0.076 \tabularnewline
\bottomrule
\end{tabularx}
\end{table}

All three void catalogues yield positive $\gamma$\textsubscript{env}, matching MTDF\_\allowbreak{}02 values to the last reported decimal place.

\subsection{NGC/\allowbreak{}SGC split}
\begin{table}[H]
\centering
\footnotesize
\begin{tabularx}{\linewidth}{l>{\RaggedRight\arraybackslash}X>{\RaggedRight\arraybackslash}X}
\toprule
\textbf{Footprint} & \textbf{$\gamma$\textsubscript{env} (REVOLVER)} & \textbf{Interpretation} \tabularnewline
NGC & +0.0050 & Strong positive signal \tabularnewline
SGC & +0.0008 & Null, consistent with reduced void coverage \tabularnewline
\bottomrule
\end{tabularx}
\end{table}

The NGC/\allowbreak{}SGC contrast is consistent with footprint-dependent detectability under a finite-range stress field. Dedicated diagnostic tests are planned (see Section 7).

\subsection{Survey fixed effects}
All three catalogues show |$\Delta$$\gamma$\textsubscript{env}| < 0.0005 when survey fixed effects are included. The environment signal is not an artefact of survey-dependent systematics.

\subsection{Leave-one-survey-out (LOSO)}
16 of 16 individual survey removals retain positive $\gamma$\textsubscript{env}. The signal is stable under jackknife.

\subsection{Redshift modulation: the signature result}
\begin{table}[H]
\centering
\footnotesize
\begin{tabularx}{\linewidth}{l>{\RaggedRight\arraybackslash}X>{\RaggedRight\arraybackslash}X>{\RaggedRight\arraybackslash}X>{\RaggedRight\arraybackslash}X}
\toprule
\textbf{Redshift bin} & \textbf{REVOLVER p} & \textbf{VIDE p} & \textbf{VoidFinder p} & \textbf{Interpretation} \tabularnewline
z $\in$ [0.02, 0.04) & \textbf{0.0035} & \textbf{0.0045} & 0.06 & Strong signal \tabularnewline
z > 0.04 & > 0.84 & > 0.84 & > 0.84 & Null \tabularnewline
\bottomrule
\end{tabularx}
\end{table}

\begin{mdframed}[linecolor=accent,linewidth=0.5pt,backgroundcolor=callout-bg,leftline=true,rightline=true,topline=true,bottomline=true,innerleftmargin=8pt,innerrightmargin=8pt,innertopmargin=6pt,innerbottommargin=6pt]
\textbf{Physical interpretation.}
The signal concentrates at low redshift (z~<~0.04) and disappears at higher redshifts. This is a specific, a priori MTDF prediction arising from the finite-range response of the stress field (strain-lever scale $\beta$~=~22.7~Mpc). The redshift decay structure is not built into any of the ad-hoc comparison models. $\Lambda$CDM does predict environment-correlated SN residuals through peculiar velocities and weak lensing, and Pantheon+ corrects for these using 2M++ reconstructions; the MTDF-specific signature is the excess amplitude, the sharp low-z transition, and its sign.
\end{mdframed}

\subsection{Non-parametric confirmation}

\subsubsection{Permutation test (10,000 shuffles)}
\begin{table}[H]
\centering
\footnotesize
\begin{tabularx}{\linewidth}{l>{\RaggedRight\arraybackslash}X}
\toprule
\textbf{Catalogue} & \textbf{p\textsubscript{perm}} \tabularnewline
REVOLVER & 0.025 \tabularnewline
VIDE & 0.046 \tabularnewline
VoidFinder & 0.062 \tabularnewline
\bottomrule
\end{tabularx}
\end{table}

\subsubsection{Block bootstrap (5,000 resamples), 95\% CI}
\begin{table}[H]
\centering
\footnotesize
\begin{tabularx}{\linewidth}{l>{\RaggedRight\arraybackslash}X>{\RaggedRight\arraybackslash}X}
\toprule
\textbf{Catalogue} & \textbf{95\% CI (mag)} & \textbf{Excludes zero?} \tabularnewline
REVOLVER & {}[+0.0016, +0.0076] & Yes \tabularnewline
VIDE & {}[+0.0010, +0.0080] & Yes \tabularnewline
\bottomrule
\end{tabularx}
\end{table}

\subsection{Numerical integrity}
GPU vs CPU crosscheck: maximum difference = 1.86~$\times$~10\textsuperscript{$-$17} (machine epsilon). Results are bit-reproducible.

\begin{mdframed}[linecolor=accent,linewidth=0.5pt,backgroundcolor=callout-bg,leftline=true,rightline=true,topline=true,bottomline=true,innerleftmargin=8pt,innerrightmargin=8pt,innertopmargin=6pt,innerbottommargin=6pt]
\textbf{Phase 3 conclusion.}
The SN~$\times$~void environment signal is reproduced by an independent reimplementation and survives non-parametric robustness checks. Three void catalogues, 16/\allowbreak{}16 LOSO surveys, permutation tests, and block bootstrap all support a positive $\gamma$\textsubscript{env} with the redshift decay structure predicted by MTDF, and bootstrap 95\% confidence intervals exclude zero. This constitutes an independent reimplementation and non-parametric robustness validation of the MTDF\_\allowbreak{}02 result on the same data. Independent confirmation in the scientific sense requires analysis of new data by external groups, which remains the priority follow-up (Phase 4 forecasts the required sample sizes).
\end{mdframed}

\section{Phase 4: Sensitivity forecasts}
\label{sec:sec5}
COMPLETE: 5$\sigma$ detection at \~{}3,400 SNe (LSST/\allowbreak{}Rubin territory)Phase 4 uses GPU-accelerated Monte Carlo simulations to forecast the sample sizes required for 3$\sigma$ and 5$\sigma$ detection of the environment signal under various systematic scenarios. All six consistency checks against Phase 3 pass (relative differences \~{}10\textsuperscript{$-$13}).

\subsection{Key structural insight}
\begin{mdframed}[linecolor=accent,linewidth=0.5pt,backgroundcolor=callout-bg,leftline=true,rightline=true,topline=true,bottomline=true,innerleftmargin=8pt,innerrightmargin=8pt,innertopmargin=6pt,innerbottommargin=6pt]
\textbf{Rank-1 systematic immunity.}
A global systematic floor of the form \( \sigma_{\rm sys}^2 \times \mathbf{1}\mathbf{1}^T \) has \emph{exactly zero} effect on \( \sigma_{\gamma} \). The intercept absorbs global systematics entirely. Only systematics correlated with \( d_{\rm signed} \) (the void proximity metric) can degrade the environment coefficient. Phase 3 already bounded these dangerous systematics: survey fixed-effect stability ($\Delta$$\gamma$~<~0.0005), LOSO (16/\allowbreak{}16 positive), and redshift modulation (signal at z~<~0.04 only).
\end{mdframed}

\textbf{The regression is statistics-limited, not systematics-limited.}

\subsection{Detection threshold table}
\begin{table}[H]
\centering
\footnotesize
\begin{tabularx}{\linewidth}{l>{\RaggedRight\arraybackslash}X>{\RaggedRight\arraybackslash}X>{\RaggedRight\arraybackslash}X>{\RaggedRight\arraybackslash}X>{\RaggedRight\arraybackslash}X}
\toprule
\textbf{Finder} & \textbf{Scenario} & \textbf{3$\sigma$ 50\%} & \textbf{3$\sigma$ 95\%} & \textbf{5$\sigma$ 50\%} & \textbf{5$\sigma$ 95\%} \tabularnewline
REVOLVER & Baseline & 1,200 & 3,000 & 3,400 & 6,600 \tabularnewline
REVOLVER & Realistic & 1,200 & 2,900 & 3,300 & 6,600 \tabularnewline
REVOLVER & Adversarial & >10k & >10k & >10k & >10k \tabularnewline
REVOLVER & NGC-only & 3,200 & 7,300 & 8,300 & >10k \tabularnewline
VoidFinder & Baseline & 1,700 & 4,500 & 4,900 & 9,100 \tabularnewline
VIDE & Baseline & 1,500 & 4,400 & 4,500 & 8,700 \tabularnewline
\bottomrule
\end{tabularx}
\end{table}

\subsection{Interpretation}
\begin{itemize}
  \item \textbf{Current Pantheon+ (564 SNe):} \~{}2$\sigma$, consistent with Phase 3 ($\Delta$$\chi$$^{2}$~=~4.25 for REVOLVER).
  \item \textbf{Baseline $\approx$ Realistic:} nearly identical thresholds confirm the analysis is statistics-limited, not systematics-limited.
  \item \textbf{Adversarial:} a hypothetical systematic perfectly correlated with \( d_{\rm signed} \) would prevent detection at any sample size, but Phase 3 already rules this out empirically.
  \item \textbf{LSST/\allowbreak{}Rubin (10,000+ low-z SNe):} well into 5$\sigma$ detection territory for all three void catalogues.
\end{itemize}

\begin{mdframed}[linecolor=accent,linewidth=0.5pt,backgroundcolor=callout-bg,leftline=true,rightline=true,topline=true,bottomline=true,innerleftmargin=8pt,innerrightmargin=8pt,innertopmargin=6pt,innerbottommargin=6pt]
\textbf{Merged-sample hardening:} A merged Pantheon+ / ZTF DR2 / Foundation DR1 sample (4,188 SNe, 3,082 at z < 0.157) confirms that the constant-$\gamma$ global model is null ($\Delta$$\chi$\textsuperscript{2} < 1.4), while a piecewise transition at z $\approx$ 0.04 yields $\Delta$$\chi$\textsuperscript{2} = 36--39 (diagonal covariance; p < 0.001) across all three void catalogues. SALT2 population controls shift $\gamma$ by < 0.3$\sigma$. The signal is localised to the finite-range regime, not a global tilt. See MTDF\_\allowbreak{}02 \S{}2.3 for full merged-sample hardening results.
\end{mdframed}

\begin{mdframed}[linecolor=accent,linewidth=0.5pt,backgroundcolor=callout-bg,leftline=true,rightline=true,topline=true,bottomline=true,innerleftmargin=8pt,innerrightmargin=8pt,innertopmargin=6pt,innerbottommargin=6pt]
\textbf{Phase 4 conclusion.}
Definitive 5$\sigma$ detection of the MTDF environment signal requires approximately 3,400 low-redshift supernovae cross-matched with void catalogues (REVOLVER baseline, 50\% power). This is within the expected yield of LSST/\allowbreak{}Rubin. The analysis is statistics-limited: realistic systematics do not materially degrade detection prospects.
\end{mdframed}

\section{5b. Phase 3b: Asymmetry and detectability diagnostics}
\label{sec:sec5b}
COMPLETE: SGC sensitivity-limited; asymmetry not significant; geometry excludedTo resolve the observed North--South Galactic Cap (NGC/\allowbreak{}SGC) asymmetry in the Phase~3 environment signal, we implemented an expanded diagnostic battery designed to distinguish physical inconsistency from detectability limits imposed by current survey coverage. All five tests ran in production mode (20s total). All diagnostics reuse the exact Phase~3 environment definition (d\textsubscript{signed} array) and regression structure. No void assignments, distances, or environment metrics are recomputed.

\subsection{5b.1 Test A: IPW matched redshift}
After inverse probability weighting to equalise footprint demographics (redshift, host mass proxy, survey origin), no catalogue produces a significant $\Delta$$\gamma$. All p-values range from 0.36 to 0.99. REVOLVER shows the largest residual at +0.0063, still far from significance. Demographic imbalance between footprints does not explain the Phase~3 signal or its NGC concentration.

\subsection{5b.2 Test B: Void--SN joint density mapping}
The SGC is severely void-sparse relative to the NGC across all three catalogues:

\begin{table}[H]
\centering
\footnotesize
\begin{tabularx}{\linewidth}{l>{\RaggedRight\arraybackslash}X>{\RaggedRight\arraybackslash}X>{\RaggedRight\arraybackslash}X}
\toprule
\textbf{Catalogue} & \textbf{NGC voids} & \textbf{SGC voids} & \textbf{Ratio} \tabularnewline
VoidFinder & 264 & 35 & 7.5$\times$ \tabularnewline
REVOLVER & 511 & 82 & 6.2$\times$ \tabularnewline
VIDE & 383 & 58 & 6.6$\times$ \tabularnewline
\bottomrule
\end{tabularx}
\end{table}

SN-per-void ratios are 3--8$\times$ higher in the SGC, indicating that each SGC void is sampled by far fewer supernovae. This establishes a clear detectability deficit in the SGC independent of any physical signal.

\subsection{5b.3 Test C: Wald test and parametric bootstrap}
The Wald test for $\Delta$$\gamma$ = $\gamma$\textsubscript{NGC} $-$ $\gamma$\textsubscript{SGC} returns p-values of 0.49--0.96 across catalogues, indicating no significant asymmetry. The parametric bootstrap (likelihood ratio test under H\textsubscript{0}: $\gamma$\textsubscript{NGC} = $\gamma$\textsubscript{SGC}) returns p = 0.17--0.24. Neither test rejects the null hypothesis that both footprints share the same underlying $\gamma$\textsubscript{env}.

\subsection{5b.4 Test D: Geometry and mode-coupling check}
Isotropic mock catalogues preserving the angular mask and redshift selection function produce $\Delta$$\gamma$ distributions fully consistent with zero. Empirical p-values for the observed $\Delta$$\gamma$ range from 0.28 to 0.90. Survey geometry alone cannot generate the observed hemispherical difference.

\subsection{5b.5 Test E: Null signal injection (SGC recovery)}
An NGC-strength $\gamma$\textsubscript{env} signal was injected into SGC supernova distance moduli using the exact Phase~3 d\textsubscript{signed} values:

\begin{table}[H]
\centering
\footnotesize
\begin{tabularx}{\linewidth}{l>{\RaggedRight\arraybackslash}X>{\RaggedRight\arraybackslash}X>{\RaggedRight\arraybackslash}X}
\toprule
\textbf{Catalogue} & \textbf{Recovery bias} & \textbf{Detection rate at 1$\times$ (2$\sigma$)} & \textbf{Detection rate at 1.5$\times$} \tabularnewline
REVOLVER & <3\% & 44\% & -- \tabularnewline
VIDE & -- & 16\% & -- \tabularnewline
VoidFinder & -- & \ {}0\% & -- \tabularnewline
\bottomrule
\end{tabularx}
\end{table}

VoidFinder's NGC $\gamma$ is near zero, so injection at 1$\times$ is negligible. REVOLVER is the only catalogue with meaningful SGC recovery power.

REVOLVER recovers the injected signal with less than 3\% bias, but achieves only 44\% detection rate at 2$\sigma$ for a 1$\times$ NGC-strength injection. This mechanically demonstrates that the SGC lacks the statistical power to detect the signal even when it is present.

\subsection{5b.6 Interpretation}
\begin{mdframed}[linecolor=accent,linewidth=0.5pt,backgroundcolor=callout-bg,leftline=true,rightline=true,topline=true,bottomline=true,innerleftmargin=8pt,innerrightmargin=8pt,innertopmargin=6pt,innerbottommargin=6pt]
\textbf{Phase 3b conclusion.}
Three findings, each supported by independent diagnostics:
\begin{enumerate}
  \item \textbf{Not significant.} The NGC/\allowbreak{}SGC contrast is not statistically significant (Wald p = 0.49--0.96; bootstrap p = 0.17--0.24). After IPW demographic matching, all $\Delta$$\gamma$ are insignificant (p = 0.36--0.99).
  \item \textbf{Not geometric.} Isotropic mocks with real mask and selection produce null $\Delta$$\gamma$ distributions consistent with zero (p = 0.28--0.90). Survey geometry cannot generate the observed pattern.
  \item \textbf{Sensitivity-limited.} The SGC has 6--8$\times$ fewer voids than the NGC. Null injection tests confirm that an NGC-strength signal is recoverable in SGC-like data only at 44\% power (REVOLVER, 2$\sigma$). The SGC null result reflects insufficient void--SN density, not physical inconsistency with MTDF.
\end{enumerate}
\end{mdframed}

\section[Phase 5: Full Planck plik + class\_mtdf (v1 retired; sealed v2)]{Phase 5: Full Planck plik + class\_\allowbreak{}mtdf (superseded by the v2 chains)}
\label{sec:sec6}
v1 RESULT RETIRED (gauge artifact). v2 SEALED: standard parameters unshifted, k\textsubscript{f} = 0.685 $\pm$ 0.500, $\sigma$\textsubscript{8} = 0.8265 $\pm$ 0.0062 (enhanced growth). $\Delta$$\chi$$^{2}$ vs matched $\Lambda$CDM control = $-$1.29 (bounded best fits). This is a maximum-likelihood comparison, not a Bayes factor or Bayesian evidence calculation.

Phase 5 was the Priority~1 hard falsifier: running the full Planck 2018 plik TTTEEE + lowl TT + lowl EE + lensing likelihood using the \texttt{class\_\allowbreak{}mtdf} modified Boltzmann solver. The v1 execution of this phase (minimisation plus converged MCMC) is retained in the archived V74 release and in \texttt{validation/\allowbreak{}output/\allowbreak{}phase5/\allowbreak{}} as a historical record. Its numerical results are \textbf{void}, for the reason documented below. They are replaced here by the sealed v2 chain results.

\begin{mdframed}[linecolor=retire,linewidth=0.5pt,backgroundcolor=callout-bg,leftline=true,rightline=true,topline=true,bottomline=true,innerleftmargin=8pt,innerrightmargin=8pt,innertopmargin=6pt,innerbottommargin=6pt]
\textbf{Why the v1 numbers are void: the gauge artifact.}
The v1 \texttt{class\_\allowbreak{}mtdf} implementation applied the growth modification $\mu$ to the 00 (density) source term in the synchronous gauge. That insertion breaks the near-cancellation between the k$^{2}$$\eta$ and 1.5\,a$^{2}$$\delta$$\rho$ terms in the synchronous-gauge Poisson bookkeeping and flips the apparent sign of the growth response. The chains returned $\sigma$\textsubscript{8} = 0.790 $\pm$ 0.006 and S\textsubscript{8} = 0.802 $\pm$ 0.013, an apparent clustering suppression read at the time as easing the S\textsubscript{8} tension. The corrected implementation applies the covariant, k-independent $\mu$(a) in the newtonian gauge (CDM Euler equation), where the modification is gauge-consistent. Under the corrected implementation the growth response is an \emph{enhancement}, as the analytic derivation requires ($\mu$ > 1 at all late redshifts). The suppression result, the S\textsubscript{8}-relief reading, the v1 posterior tables, and the v1 robustness numbers built on those chains (k\textsubscript{f} = 0.495/\allowbreak{}0.50 $\pm$ 0.36, $\Delta$$\chi$$^{2}$ = +0.63, $\Delta$AIC = +2.63, $\Delta$BIC = +8.05) are all downstream of the artifact. They carry no evidential weight.

\textbf{Permanent methods item (gauge-invariance regression test).} Any modified-gravity source term inserted into a Boltzmann solver must be validated by a gauge-invariance regression. Implement the modification independently in synchronous and newtonian gauges, and require agreement of the gauge-invariant growth observables (f$\sigma$\textsubscript{8}, $\sigma$\textsubscript{8}(z), C\textsubscript{$\ell$}) before any chain is run. This test is now a permanent part of the MTDF validation methodology. It is exactly the test that would have caught the v1 artifact before publication.
\end{mdframed}

\subsection{Sealed v2 results (corrected implementation)}

\begin{mdframed}[linecolor=accent,linewidth=0.5pt,backgroundcolor=callout-bg,leftline=true,rightline=true,topline=true,bottomline=true,innerleftmargin=8pt,innerrightmargin=8pt,innertopmargin=6pt,innerbottommargin=6pt]
\textbf{Apples-to-apples guarantee (unchanged).}
Both $\Lambda$CDM and MTDF configurations use identical likelihood set (plik TTTEEE + lowl TT + lowl EE + lensing), data bins, masks, nuisance list, nuisance priors, cosmological priors, and stop rules. The only differences are the added free parameter k\textsubscript{f} and the corrected \texttt{class\_\allowbreak{}mtdf} solver in place of standard CLASS.
\end{mdframed}

\textbf{Run:} four chains, \~{}66 h, frozen flight configuration (hash-pinned binary and YAML). The chains stopped at the frozen hard cap of 200,000 accepted steps per chain, with R$-$1 = 0.0085 (target 0.01, satisfied). R$-$1\textsubscript{cl} = 0.060 (strict 0.02 not reached: 95\% bounds carry extra sampling noise; means and peaks solid). All interruptions auto-resumed from checkpoint with no sample loss.

\begin{table}[H]
\centering
\footnotesize
\begin{tabularx}{\linewidth}{l>{\RaggedRight\arraybackslash}X>{\RaggedRight\arraybackslash}X>{\RaggedRight\arraybackslash}X}
\toprule
\textbf{Parameter} & \textbf{$\Lambda$CDM baseline} & \textbf{MTDF v2} & \textbf{$\Delta$} \tabularnewline
$\sigma$\textsubscript{8} & 0.8101 & 0.82651 $\pm$ 0.00623 & +0.0164 (+1.91$\sigma$, posterior units) \tabularnewline
H\textsubscript{0} (km/\allowbreak{}s/\allowbreak{}Mpc) & 67.382 & 67.444 $\pm$ 0.546 & +0.06 \tabularnewline
$\Omega$\textsubscript{m} & -- & 0.31533 $\pm$ 0.00771 & negligible \tabularnewline
n\textsubscript{s} & -- & 0.96426 & negligible \tabularnewline
ln(10\textsuperscript{10}A\textsubscript{s}) & -- & 3.03962 & negligible \tabularnewline
k\textsubscript{f} & -- & 0.685 $\pm$ 0.500 & 0 and 1 both compatible \tabularnewline
\bottomrule
\end{tabularx}
\end{table}

Sealed v2 reporting set (k\textsubscript{f} free), adjudicated by the sealed analysis script under a pre-registered three-posture rule. The adjudication was executed only on the completed chain. The k\textsubscript{f} = 1 importance-reweighted rider gives the same posture: $\sigma$\textsubscript{8} shift +0.0173 (+2.07$\sigma$), f\textsubscript{kick} = 0.0033, $\mu$(z\textsubscript{rec}) = 1.00002 (CMB-inert). Native growth-sector units: S\textsubscript{8} = 0.8473 $\pm$ 0.0138 (free), 0.8501 $\pm$ 0.0132 (k\textsubscript{f} = 1).

\subsection{Interpretation}
The corrected covariant $\mu$(a) enhances late growth by about +2 percent, and the full-likelihood chain confirms this end to end. $\sigma$\textsubscript{8} lands about 2$\sigma$ (posterior units) above the $\Lambda$CDM baseline. The corrected theory therefore mildly \emph{worsens} the comparison with published weak-lensing clustering amplitudes (sealed posture: S8-WORSENED, both riders agreeing). Standard-parameter shifts are tiny and the CMB-side H\textsubscript{0} is unchanged. What survives unchanged from the v1 chains is the structural finding that Planck alone cannot pin the EFE amplitude. The k\textsubscript{f} posterior is broad, with both k\textsubscript{f} = 0 and the theory value k\textsubscript{f} = 1 compatible. Planck establishes consistency, not support. Discrimination lives in low-redshift probes.

\subsection{Model comparison against the matched $\Lambda$CDM control}
The matched $\Lambda$CDM control chain uses the identical binary, likelihood set, priors, chain profile, and stop rules, with exactly three configuration deltas removing the MTDF theory keys. It reproduces the reference Planck-$\Lambda$CDM posterior at full-posterior level, validating the pipeline end to end. Control values: $\sigma$\textsubscript{8} = 0.81036 $\pm$ 0.00598, per-sample S\textsubscript{8} = 0.83004 $\pm$ 0.01278, H\textsubscript{0} = 67.389 $\pm$ 0.544, with shared-parameter shifts at the 10\textsuperscript{$-$4} level. At the bounded best fits, MTDF $-$log\,L = 1385.78 (at k\textsubscript{f} = 0.476) against $\Lambda$CDM $-$log\,L = 1386.42, giving \textbf{$\Delta$$\chi$$^{2}$ (MTDF $-$ $\Lambda$CDM) = $-$1.29}. The minimisation was seeded from each chain's best sampled point, and both minimisations converged. Dataset-resolved, the two models agree to about one point on every component (lowl TT 22.29 vs 23.77, lowl EE 396.47 vs 396.27, plik TTTEEE 2343.80 vs 2344.00, lensing 8.99 vs 8.80). The small MTDF gain sits in lowl TT, the late-ISW region. Per-sample growth-sector amplitudes: MTDF S\textsubscript{8} = 0.8473 $\pm$ 0.0138 (free) and 0.8501 $\pm$ 0.0132 (k\textsubscript{f} = 1) against the control's 0.83004 $\pm$ 0.01278.

\emph{This is a maximum-likelihood comparison, not a Bayes factor or Bayesian evidence calculation.} With one additional free parameter, a one-to-two-point gain is the no-information expectation ($\Delta$AIC = +0.71, favouring neither model). A BIC-style penalty depends on an adopted effective Planck data count and is not quoted. The full Planck 2018 likelihood does not distinguish the corrected MTDF from $\Lambda$CDM. Chain minima as reported by the sampler are minus-log-posterior values and are not compared across models. All numbers above are likelihood-only, from the bounded minima and per-dataset components.

\subsection{Scope statement}
The v2 run is a full Planck 2018 likelihood test of the corrected MTDF effective cosmological implementation. That implementation is the corrected covariant $\mu$(a)/\allowbreak{}EFE sector, with the high-Q carrier represented by its CDM-like limit (standard pressureless CDM as proxy). It is not a Boltzmann implementation of all nine degrees of freedom of the frozen action. The full-action CMB perturbation implementation (unified scalar + vector perturbations) is a named future calculation. Phase-5-level success discharges the CMB confrontation at the effective level only.

\section{6b. Phase 6: Cross-probe discriminator tests}
\label{sec:sec6b}
COMPLETE: Two discriminators (A, E), three upper limits (B, B2, D), one internal consistency (C), two forecasts (C3, C4), one resolved null (C5), one retired-as-run comparison (C2, v1-posterior-based)  -  10 tests totalPhase~6 extends validation beyond CMB compatibility toward tests that can \emph{distinguish} MTDF from $\Lambda$CDM using independent probes. Ten tests are packaged under \texttt{validation/\allowbreak{}output/\allowbreak{}phase6/\allowbreak{}}, targeting weak lensing, CMB lensing, derived parameter consistency, cross-probe S\textsubscript{8} coherence, growth--environment coupling, BAO residuals, cluster mass ratios, and the SN~$\times$~void redshift transition. All code was written from scratch for Phase~6.

Tests~A and E positively support MTDF-style predictions (redshift transition, void-specific CMB lensing signal). Test~C2, as run against the v1 posterior, is retired (gauge artifact); under the sealed v2 chains the S\textsubscript{8} comparison mildly worsens, and its quantitative revisiting is deferred to the preregistered FWD-S8 programme. Tests~B, B2, and D are systematics-clean nulls with explicit upper limits. Test~C confirms internal consistency. Tests~C3 and C4 quantify sensitivity boundaries. Test~C5 resolves the cluster mass ratio ($\eta$~=~1). Together they demonstrate: \emph{the pipelines contain no favourable selection and are sound, no hidden contradictions exist, and the boundaries of current discriminating power are mapped.}

\subsection{Summary table}
\begin{table}[H]
\centering
\footnotesize
\begin{tabularx}{\linewidth}{l>{\RaggedRight\arraybackslash}X>{\RaggedRight\arraybackslash}X>{\RaggedRight\arraybackslash}X>{\RaggedRight\arraybackslash}X}
\toprule
\textbf{Test} & \textbf{Observable} & \textbf{Primary statistic} & \textbf{Result} & \textbf{Implication} \tabularnewline
\midrule
\textbf{A} & SN $\times$ void $\gamma$\textsubscript{env} vs z\textsubscript{cut} & GLS z-score peak over z\textsubscript{cut} scan & Peak at z\textsubscript{cut}=0.030: 3.62$\sigma$ (global p < 0.005) & Supports MTDF transition at z < 0.04 \tabularnewline
\textbf{B} & KiDS-1000 $\gamma$\textsubscript{t} around void centres & $\Delta$$\gamma$\textsubscript{t} over R=5--20 Mpc/\allowbreak{}h & $-$3.65e$-$5 $\pm$ 4.68e$-$5 (0.8$\sigma$); 95\%: |$\Delta$$\gamma$\textsubscript{t}| < 1.2e$-$4 & Null  -  bounded, not yet discriminating \tabularnewline
\textbf{B2} & KiDS-1000 trough vs ridge $\gamma$\textsubscript{t} & $\Delta$$\gamma$\textsubscript{t,split} (ridge$-$trough) & +2.56e$-$5 $\pm$ 5.46e$-$5 (0.5$\sigma$); 95\%: |$\Delta$$\gamma$\textsubscript{t,split}| < 1.3e$-$4 & Null  -  bounded, not yet discriminating \tabularnewline
\textbf{C} & Derived parameter consistency & Max shift vs $\Lambda$CDM & All within 2.4$\sigma$ & Internally consistent (v1-posterior-based; the S\textsubscript{8} element is superseded) \tabularnewline
\textbf{D} & Planck PR4 $\kappa$ $\times$ DESI BGS voids & Per-void compensated $\kappa$\textsubscript{comp} & $-$7.47e$-$4 $\pm$ 1.03e$-$3 (0.7$\sigma$, 557 voids) & Null at low z  -  expected; pipeline validated \tabularnewline
\textbf{E} & Planck PR4 $\kappa$ $\times$ BOSS DR12 voids & Conservative $\kappa$\textsubscript{comp} + detection-optimised CTH/\allowbreak{}MF & Comp: 0.2$\sigma$ null; CTH: 1.5$\sigma$ (p\textsubscript{rand}=0.03); MF: 1.9$\sigma$ (p\textsubscript{rand}=0.02) & Pipeline benchmark validated; void-specific signal recovered \tabularnewline
\textbf{C2} & S\textsubscript{8} cross-probe comparison (retired as run) & Phase 5 S\textsubscript{8} vs published WL surveys & Combined WL tension: 3.47$\sigma$ $\rightarrow$ 1.89$\sigma$ (46\% reduction) & Retired as run; under the sealed v2 chains the direction reverses \tabularnewline
\textbf{C3} & BAO residual structure & $\Delta$W RMS (wiggle difference MTDF$-$$\Lambda$CDM) & $\Delta$W RMS = 0.27\%, max = 0.51\%; sound horizon shift +0.11\% & Not detectable by any foreseeable survey (best SNR = 0.92) \tabularnewline
\textbf{C4} & f$\sigma$\textsubscript{8} $\times$ environment & RSD + void f$\sigma$\textsubscript{8} vs $\mu$(a) prediction & $\Delta$$\chi$$^{2}$ = +0.49 (9 DOF); max discrimination 0.22$\sigma$ & Sensitivity-limited; needs \ {}10$\times$ improvement (DESI/\allowbreak{}Euclid) \tabularnewline
\textbf{C5} & Cluster M\textsubscript{dyn}/M\textsubscript{lens} & Analytical $\mu$(a) prediction + numerical $\eta$(k,z) verification & $\eta$ = 1.000000 at all cluster scales; Case B confirmed: M\textsubscript{dyn}/M\textsubscript{lens} = 1 & Resolved: cluster mass ratio is not a discriminator for MTDF \tabularnewline
\bottomrule
\end{tabularx}
\end{table}

\subsection{6b.1 Test A: Redshift transition scan}
\begin{itemize}
  \item Scans z\textsubscript{cut} $\in$ [0.025,~0.100] using GLS and Spearman metrics on the Phase~3 SN~$\times$~void environment signal.
  \item Peak at z\textsubscript{cut}~=~0.030 with GLS z-score 3.62$\sigma$; global scan-corrected permutation p~<~0.005 (limited by 200 permutations; 0 of 200 exceedances). The scan peak (z\textsubscript{cut}~=~0.030) and the headline transition (z~$\approx$~0.040~$\pm$~0.005) are statistically consistent given the scan's redshift resolution; they refer to the same low-redshift transition, not two distinct features.
  \item Spearman rank test agrees independently (3.1$\sigma$).
  \item Interpretation: the environment signal is localised to z~<~0.04, consistent with the Phase~3 redshift modulation and the MTDF finite-range prediction.
  \item Artefacts: validation/\allowbreak{}output/\allowbreak{}phase6/\allowbreak{}testA\_\allowbreak{}redshift\_\allowbreak{}transition/\allowbreak{}
\end{itemize}

\subsection{6b.2 Test B: Weak lensing $\times$ environment (void centres)}
\begin{itemize}
  \item KiDS-1000 shear stacked around DESI REVOLVER void centres and matched random positions; $\gamma$\textsubscript{x} (cross-shear) gate required to pass before reporting $\gamma$\textsubscript{t}.
  \item $\Delta$$\gamma$\textsubscript{t}~=~$-$3.65$\times$10\textsuperscript{$-$5}~$\pm$~4.68$\times$10\textsuperscript{$-$5} over R~=~5--20~Mpc/\allowbreak{}h, consistent with zero.
  \item $\gamma$\textsubscript{x} gate: pass (p~=~0.234).
  \item Explicit two-sided 95\% upper limit: |$\Delta$$\gamma$\textsubscript{t}|~<~1.20$\times$10\textsuperscript{$-$4}.
  \item Without a forward MTDF lensing prediction for this selection, the result is best framed as ``not yet discriminating, but now bounded and reproducible.''
  \item Artefacts: validation/\allowbreak{}output/\allowbreak{}phase6/\allowbreak{}testB\_\allowbreak{}wl\_\allowbreak{}environment/\allowbreak{}
\end{itemize}

\subsection{6b.3 Test B2: Trough and ridge lensing (projected density split)}
\begin{itemize}
  \item Projected density field built with mask-aware normalisation and full KiDS-1000 footprint (North + South). Both trough and ridge $\gamma$\textsubscript{x} gates pass.
  \item $\Delta$$\gamma$\textsubscript{t,split} (ridge$-$trough)~=~+2.56$\times$10\textsuperscript{$-$5}~$\pm$~5.46$\times$10\textsuperscript{$-$5} (0.5$\sigma$), consistent with zero.
  \item Explicit two-sided 95\% upper limit: |$\Delta$$\gamma$\textsubscript{t,split}|~<~1.28$\times$10\textsuperscript{$-$4}.
  \item Projected P20/\allowbreak{}P80 density contrasts ($\delta$~$\sim$~$\pm$0.15) are too mild for detectable differential lensing at KiDS sensitivity. Spectroscopic 3D density or LSST/\allowbreak{}DES area would be needed for significance.
  \item Artefacts: validation/\allowbreak{}output/\allowbreak{}phase6/\allowbreak{}testB2\_\allowbreak{}trough\_\allowbreak{}ridge/\allowbreak{}
\end{itemize}

\subsection{6b.4 Test C: Derived parameter consistency}
\begin{itemize}
  \item Cross-checks derived parameters from the Phase~5 MCMC posterior against external weak lensing constraints.
  \item No derived parameter deviates beyond 2.4$\sigma$ from $\Lambda$CDM.
  \item MTDF's Planck posterior is internally consistent; no hidden contradictions introduced by \texttt{class\_\allowbreak{}mtdf}.
  \item Artefacts: validation/\allowbreak{}output/\allowbreak{}phase6/\allowbreak{}testC\_\allowbreak{}derived\_\allowbreak{}consistency/\allowbreak{}
\end{itemize}

\subsection{6b.5 Test D: CMB lensing $\times$ DESI BGS voids (low redshift)}
\begin{itemize}
  \item Planck PR4 MV convergence ($\kappa$) stacked on DESIVAST BGS voids (z~<~0.24) with mask-weighted stacking. Primary statistic is a per-void compensated filter: $\kappa$(R/\allowbreak{}R\textsubscript{v}<0.5)~$-$~$\kappa$(4<R/\allowbreak{}R\textsubscript{v}<5), designed to remove large-scale mode contamination.
  \item $\kappa$\textsubscript{comp}~=~$-$7.47$\times$10\textsuperscript{$-$4}~$\pm$~1.03$\times$10\textsuperscript{$-$3} (0.7$\sigma$, 557 voids), consistent with zero.
  \item RA-scramble null: p~=~0.425 (compensated). Random-sky null: p~=~0.080.
  \item Robustness: stable across outer annulus choices (all <1.1$\sigma$), low-l removal (l$\ge$20, l$\ge$30), and void finder.
  \item Null at z~<~0.24 is expected: the CMB lensing kernel peaks at z~\~{}~1--2, and the BGS void count (557 usable) limits statistical power.
  \item Diagnosed failure modes from v1: hard mask quality cut rejected 73\% of voids; low-l CMB modes created a positive $\kappa$ pedestal. Both fixed in v2 via per-void compensated filter and mask-weighted stacking.
  \item Artefacts: validation/\allowbreak{}output/\allowbreak{}phase6/\allowbreak{}testD\_\allowbreak{}cmb\_\allowbreak{}lensing/\allowbreak{}
\end{itemize}

\subsection{6b.6 Test E: CMB lensing $\times$ BOSS DR12 voids (benchmark)}
\begin{itemize}
  \item Benchmark run on the Mao+2017 ZOBOV catalogue (1,228 BOSS DR12 CMASS+LOWZ voids, z~=~0.2--0.7) to validate the CMB lensing stacking pipeline against Cai+2017's published 3.2$\sigma$ detection on the same catalogue.
  \item Five statistics computed side by side:
    
\textbf{Per-void compensated} (systematics-clean primary): 0.2$\sigma$, null by design.
\textbf{AP disc} (contamination monitor): 4.7$\sigma$, but RA-scramble p~=~0.945 and collapses to 0.3$\sigma$ under l$\ge$20 cut  -  entirely pedestal.
\textbf{CTH} (Cai-class compensated top-hat): 1.5$\sigma$ (p\textsubscript{rand}~=~0.03), increases to 2.0$\sigma$ under l$\ge$20 cut.
\textbf{Matched filter} (template + jackknife covariance): 1.9$\sigma$ (p\textsubscript{rand}~=~0.02), increases to 2.2$\sigma$ under l$\ge$20 cut.
\textbf{Close comp} (alternative aperture): 0.5$\sigma$, consistent with zero.
  \item Detection-optimised statistics (CTH and MF) show the expected hierarchy and respond correctly to pedestal removal and footprint-preserving null tests. The gap from \~{}1.5--2.2$\sigma$ to Cai+2017's 3.2$\sigma$ reflects template shape (step vs theory-derived) and different Planck products (PR4 vs PR2015).
  \item Artefacts: validation/\allowbreak{}output/\allowbreak{}phase6/\allowbreak{}testE\_\allowbreak{}boss\_\allowbreak{}benchmark/\allowbreak{}
\end{itemize}

\subsection{6b.7 Test C2: S\textsubscript{8} cross-probe comparison (RETIRED as run; direction reversed under v2)}
\begin{itemize}
  \item \textbf{As-run status: RETIRED.} Test C2 compared the v1 Phase-5 posterior (S\textsubscript{8} = 0.802 $\pm$ 0.013) against published weak-lensing constraints and reported a tension reduction. The v1 posterior was a gauge artifact (Section 6), so the tension-reduction result is void.
  \item \textbf{Under the sealed v2 chains the direction reverses:} the corrected implementation gives an enhanced-growth posture ($\sigma$\textsubscript{8} = 0.8265 $\pm$ 0.0062, about +1.9$\sigma$ posterior units above the $\Lambda$CDM baseline; native S\textsubscript{8} = 0.8473 $\pm$ 0.0138), which mildly \emph{worsens} the comparison with published weak-lensing compressions. No relief claim survives.
  \item The as-run caveat stands and is now load-bearing: published WL S\textsubscript{8} values assume $\Lambda$CDM growth and lensing kernels, so the only pre-registered path for quantitatively revisiting this comparison is the forward-modelling programme (FWD-S8), which forward-models each survey's compression within MTDF rather than comparing native posteriors to $\Lambda$CDM-compressed numbers.
  \item Artefacts (v1, historical): validation/\allowbreak{}output/\allowbreak{}phase6/\allowbreak{}testC2\_\allowbreak{}s8\_\allowbreak{}coherence/\allowbreak{}
\end{itemize}

\subsection{6b.8 Test C3: BAO residual structure}
\begin{itemize}
  \item Extracts BAO wiggles W(k)~=~P(k)/P\textsubscript{nw}(k)~$-$~1 from \texttt{class\_\allowbreak{}mtdf} power spectra using an Eisenstein~\&~Hu (1998) no-wiggle reference, and computes the residual $\Delta$W~=~W\textsubscript{MTDF}~$-$~W\textsubscript{$\Lambda$CDM}.
  \item Sound horizon shift: +0.106\% (147.148~$\rightarrow$~147.303~Mpc). Wiggle difference: $\Delta$W RMS~=~0.27\%, max~=~0.51\% at k~=~0.12~h/\allowbreak{}Mpc. No detectable peak phase shift at grid resolution.
  \item Best projected SNR: 0.92 (DESI + Euclid combined). The 0.1\% sound horizon shift is absorbed into standard cosmological parameter degeneracies.
  \item Conclusion: BAO residual structure is not a viable MTDF--$\Lambda$CDM discriminator with any foreseeable survey.
  \item Artefacts: validation/\allowbreak{}output/\allowbreak{}phase6/\allowbreak{}testC3\_\allowbreak{}bao\_\allowbreak{}residuals/\allowbreak{}
\end{itemize}

\subsection{6b.9 Test C4: f$\sigma$\textsubscript{8} $\times$ environment}
\begin{itemize}
  \item Compares MTDF and $\Lambda$CDM f$\sigma$\textsubscript{8}(z) predictions against a compilation of 8 RSD surveys plus void-specific growth measurements from Hamaus+2020.
  \item $\chi$$^{2}$: $\Lambda$CDM~=~7.61, MTDF~=~8.10 ($\Delta$$\chi$$^{2}$~=~+0.49, 9~DOF). Maximum per-point discrimination: 0.22$\sigma$.
  \item The homogeneous $\mu$(a) modification produces a \~{}1.5\% signal, which is 5--10$\times$ smaller than current RSD measurement errors (0.03--0.10 per z-bin).
  \item Future requirement: 2$\sigma$ discrimination needs $\sigma$(f$\sigma$\textsubscript{8})~\~{}~0.003--0.004 per z-bin, achievable with DESI/\allowbreak{}Euclid-era RSD or environment-resolved estimators.
  \item Artefacts: validation/\allowbreak{}output/\allowbreak{}phase6/\allowbreak{}testC4\_\allowbreak{}fsigma8\_\allowbreak{}environment/\allowbreak{}
\end{itemize}

\subsection{6b.10 Test C5/\allowbreak{}C5b: Cluster M\textsubscript{dyn}/M\textsubscript{lens}  -  Resolved}
\begin{itemize}
  \item C5 computed the predicted ratio M\textsubscript{dyn}/M\textsubscript{lens} under two bracketing coupling cases:
    
\textbf{Case~A} ($\mu$-only dynamics, $\Sigma$~=~1): M\textsubscript{dyn}/M\textsubscript{lens}~=~1.053 at z~=~0 (+5.3\% offset).
\textbf{Case~B} (equal coupling, $\Sigma$~=~$\mu$): M\textsubscript{dyn}/M\textsubscript{lens}~=~1.000 (no offset).
  \item \textbf{C5b resolved the key unknown} by extracting $\eta$(k,z)~=~$\Phi$(k,z)/$\Psi$(k,z) directly from \texttt{class\_\allowbreak{}mtdf} Newtonian-gauge transfer functions across k~=~0.005--2~h/\allowbreak{}Mpc and z~=~0--10.
  \item Result: \textbf{$\eta$~=~1.000000} at all cluster-relevant scales (z~$\le$~0.8). Max~|$\Delta$$\eta$| between MTDF and $\Lambda$CDM:~0.000000. Tiny deviations (|$\eta$$-$1|~\~{}~10\textsuperscript{$-$5}) appear only at z~$\ge$~5 from standard radiation anisotropic stress, present equally in both models.
  \item Physical reason: \texttt{class\_\allowbreak{}mtdf} modifies only the Poisson equation source ($\delta$$\rho$~$\rightarrow$~$\mu$$\cdot$$\delta$$\rho$) and the CDM Euler equation. The $\Phi$$-$$\Psi$ anisotropy equation is unchanged from GR. Both potentials are enhanced equally $\rightarrow$ $\Sigma$~=~$\mu$ $\rightarrow$ \textbf{Case~B confirmed}.
  \item Implication: M\textsubscript{dyn}/M\textsubscript{lens}~=~1.000 under MTDF. The cluster mass ratio channel is \textbf{not a discriminator}. Clusters can still test MTDF through growth-dependent observables (e.g.~cluster counts), because $\mu$~>~1 modifies the growth rate.
  \item Artefacts: validation/\allowbreak{}output/\allowbreak{}phase6/\allowbreak{}testC5\_\allowbreak{}cluster\_\allowbreak{}mass/\allowbreak{}, validation/\allowbreak{}output/\allowbreak{}phase6/\allowbreak{}testC5b\_\allowbreak{}gravitational\_\allowbreak{}slip/\allowbreak{}
\end{itemize}

\begin{mdframed}[linecolor=accent,linewidth=0.5pt,backgroundcolor=callout-bg,leftline=true,rightline=true,topline=true,bottomline=true,innerleftmargin=8pt,innerrightmargin=8pt,innertopmargin=6pt,innerbottommargin=6pt]
\textbf{Phase 6 conclusion (10 tests, all complete).}
The ten Phase~6 tests fall into five categories. \emph{Discriminators} (6A, 6E): the redshift transition is identified at 3.6$\sigma$; the BOSS CMB lensing benchmark recovers void-specific structure at 1.5--1.9$\sigma$. \emph{Coherence} (6C, 6C2): derived parameters are self-consistent (no shift > 2.4$\sigma$); the as-run 6C2 S\textsubscript{8} comparison is retired with the v1 chains (direction reversed under v2). \emph{Upper limits} (6B, 6B2, 6D): systematics-clean nulls with explicit 95\% bounds on void lensing and CMB lensing signals. \emph{Forecasts} (6C3, 6C4): BAO wiggles and f$\sigma$\textsubscript{8} quantify what future surveys need to reach. \emph{Resolved} (6C5): cluster M\textsubscript{dyn}/M\textsubscript{lens} is null ($\eta$~=~1). No surviving test contradicts MTDF at the consistency level. The strongest positive outcomes are the discriminators (6A redshift transition, 6E void-specific CMB lensing structure); the former 6C2 relief claim does not survive the v2 correction and is retired.
\end{mdframed}

\section{Synthesis and framework narrative}
\label{sec:sec7}

\subsection{Summary of results (17/\allowbreak{}17 phases complete)}
\begingroup\scriptsize
\begin{xltabular}{\linewidth}{>{\RaggedRight\arraybackslash}X>{\RaggedRight\arraybackslash}X>{\RaggedRight\arraybackslash}X>{\RaggedRight\arraybackslash}X>{\RaggedRight\arraybackslash}X}
\toprule
\endfirsthead
\toprule\endhead
\bottomrule\endfoot
\textbf{Phase} & \textbf{Type} & \textbf{Status} & \textbf{Key result} & \textbf{Artefacts} \tabularnewline
Phase 1 & Discriminator & PASSED & V18 strict totals reproduced to $\delta$ = 0.000049 & validation/\allowbreak{}output/\allowbreak{}phase1/\allowbreak{} \tabularnewline
Phase 2 & Discriminator & PASSED & CMB plik-lite compatible ($\Delta$$\chi$$^{2}$ < 1); k\textsubscript{f} = 1 within 95\% in all MCMC runs & validation/\allowbreak{}output/\allowbreak{}phase2/\allowbreak{} \tabularnewline
Phase 3 & Discriminator & PASSED & SN $\times$ void signal reproduced exactly; bootstrap excludes zero & validation/\allowbreak{}output/\allowbreak{}phase3/\allowbreak{} \tabularnewline
Phase 3b & Upper limit & COMPLETE & NGC/\allowbreak{}SGC asymmetry: not significant, not geometric, sensitivity-limited (5 diagnostics) & validation/\allowbreak{}output/\allowbreak{}phase3b/\allowbreak{} \tabularnewline
Phase 4 & Forecast & COMPLETE & 5$\sigma$ at \ {}3,400 SNe (LSST/\allowbreak{}Rubin); statistics-limited & validation/\allowbreak{}output/\allowbreak{}phase4/\allowbreak{} \tabularnewline
Phase 5 (v1) & Discriminator & RETIRED & v1 minimisation and MCMC results void (gauge artifact, Section 6); archived as historical record. & \multirow{2}{=}{validation/\allowbreak{}output/\allowbreak{}phase5/\allowbreak{}} \tabularnewline
Phase 5 (v2, sealed) & Discriminator & PASSED (consistency) & Corrected implementation, R$-$1 = 0.0085. $\sigma$\textsubscript{8} = 0.8265 $\pm$ 0.0062 (enhanced growth, +1.91$\sigma$ posterior units above $\Lambda$CDM baseline); k\textsubscript{f} = 0.685 $\pm$ 0.500, both 0 and 1 compatible. $\Delta$$\chi$$^{2}$ vs matched control complete (see Section 6.3). &  \tabularnewline
Phase 6A & Discriminator & PASSED & Redshift transition scan: signal peaks at z < 0.03 (3.6$\sigma$ GLS, global p < 0.005). Spearman confirms (3.1$\sigma$). & \multirow{10}{=}{validation/\allowbreak{}output/\allowbreak{}phase6/\allowbreak{}} \tabularnewline
Phase 6B & Upper limit & COMPLETE & Weak lensing $\times$ void environment (KiDS-1000): $\Delta$$\gamma$\textsubscript{t} = $-$3.65e$-$5 $\pm$ 4.68e$-$5 (0.8$\sigma$). $\gamma$\textsubscript{x} gate pass. 95\% UL: |$\Delta$$\gamma$\textsubscript{t}| < 1.2e$-$4. &  \tabularnewline
Phase 6B2 & Upper limit & COMPLETE & Trough/\allowbreak{}ridge lensing (KiDS-1000): $\Delta$$\gamma$\textsubscript{t,split} = +2.56e$-$5 $\pm$ 5.46e$-$5 (0.5$\sigma$). Both $\gamma$\textsubscript{x} gates pass. 95\% UL: |$\Delta$$\gamma$\textsubscript{t,split}| < 1.3e$-$4. &  \tabularnewline
Phase 6C & Coherence & PASSED & Derived parameter consistency: no shifts > 2.4$\sigma$ (v1-posterior-based; the S\textsubscript{8} element is superseded, see 6b.7). &  \tabularnewline
Phase 6C2 & Coherence & RETIRED (as run) & Compared the v1 Phase-5 posterior against WL surveys; the v1 posterior is void (gauge artifact), so the reported tension reduction is retired. Direction reversed under sealed v2; quantitative revisit deferred to FWD-S8. &  \tabularnewline
Phase 6C3 & Forecast & COMPLETE & BAO residual structure: $\Delta$W RMS = 0.27\%, sound horizon shift +0.11\%. Not detectable by any current or projected survey (best SNR = 0.92). &  \tabularnewline
Phase 6C4 & Forecast & COMPLETE & f$\sigma$\textsubscript{8} $\times$ environment: RSD + void growth vs $\mu$(a) prediction. Max 0.22$\sigma$; needs DESI/\allowbreak{}Euclid-era precision (\ {}10$\times$ improvement). &  \tabularnewline
Phase 6C5/\allowbreak{}C5b & Closed null & RESOLVED & Cluster M\textsubscript{dyn}/M\textsubscript{lens}: $\eta$ = 1 confirmed numerically (C5b) $\Rightarrow$ $\Sigma$ = $\mu$, so M\textsubscript{dyn}/M\textsubscript{lens} = 1 exactly. Mass ratio offset is zero in the current MTDF perturbation model. &  \tabularnewline
Phase 6D & Upper limit & COMPLETE & CMB lensing $\times$ DESI BGS voids: $\kappa$\textsubscript{comp} = $-$7.47e$-$4 $\pm$ 1.03e$-$3 (0.7$\sigma$, 557 voids). RA-scramble clean. Null expected at z < 0.24. &  \tabularnewline
Phase 6E & Discriminator & COMPLETE & CMB lensing $\times$ BOSS DR12 benchmark: CTH 1.5$\sigma$ (p\textsubscript{rand}=0.03), MF 1.9$\sigma$ (p\textsubscript{rand}=0.02). Pipeline validated against Cai+2017. &  \tabularnewline
\end{xltabular}
\endgroup

\subsection{Claim status taxonomy}
Each MTDF claim or test outcome is assigned one of four status labels. This taxonomy is intended to provide referees and readers with an unambiguous assessment of where the framework stands.

\begin{table}[H]
\centering
\footnotesize
\begin{tabularx}{\linewidth}{l>{\RaggedRight\arraybackslash}X}
\toprule
\textbf{Label} & \textbf{Definition} \tabularnewline
\checkmark{} Validated & Independently reproduced with quantitative agreement. Robust to systematics tests. \tabularnewline
$\blacksquare$ Constrained & Consistent with data but not uniquely favoured. Current data cannot distinguish from $\Lambda$CDM. \tabularnewline
$\blacktriangle$ Hypothesis & Theoretically motivated but not yet tested against data, or tested only indirectly. \tabularnewline
$\times$ Falsifier & Identified test that could refute the framework. Not yet executed. \tabularnewline
\bottomrule
\end{tabularx}
\end{table}

\subsubsection{Status by claim}
\begingroup\scriptsize
\begin{xltabular}{\linewidth}{>{\RaggedRight\arraybackslash}X>{\RaggedRight\arraybackslash}X>{\RaggedRight\arraybackslash}X}
\toprule
\endfirsthead
\toprule\endhead
\bottomrule\endfoot
\textbf{Claim / element} & \textbf{Status} & \textbf{Evidence} \tabularnewline
V18 strict $\chi$$^{2}$ totals & \checkmark{} Validated & Phase 1: independent reproduction to $\delta$ = 0.000049 \tabularnewline
Pantheon+ SN likelihood & \checkmark{} Validated & Phase 1: $\chi$$^{2}$ = 2012.72 matched exactly \tabularnewline
DESI Y1 BAO fit & \checkmark{} Validated & Phase 1: $\chi$$^{2}$ = 15.39 matched exactly \tabularnewline
Cosmic chronometers H(z) & \checkmark{} Validated & Phase 1: $\chi$$^{2}$ = 6.98 matched exactly \tabularnewline
DR16 f$\sigma$\textsubscript{8} growth & \checkmark{} Validated & Phase 1: $\chi$$^{2}$ = 2.00 matched exactly \tabularnewline
CMB compatibility (plik-lite) & $\blacksquare$ Constrained & Phase 2: $\Delta$$\chi$$^{2}$ < 1; k\textsubscript{f} = 1 within 95\% CI. Perturbative correction layer, not full Boltzmann. \tabularnewline
CMB compatibility (full plik + class\_\allowbreak{}mtdf) & \checkmark{} Validated & Sealed v2 chains (corrected implementation): standard parameters essentially unshifted; k\textsubscript{f} = 0.685 $\pm$ 0.500 with 0 and 1 both compatible; apples-to-apples configuration; $\Delta$$\chi$$^{2}$ vs matched $\Lambda$CDM control complete (see Section 6.3). (v1 Phase-5 numbers retired, gauge artifact.) \tabularnewline
SN $\times$ void environment signal ($\gamma$\textsubscript{env}) & \checkmark{} Validated & Phase 3: exact reproduction, 3 void finders, bootstrap excludes zero \tabularnewline
Redshift decay structure (z < 0.04) & \checkmark{} Validated & Phase 3: a priori prediction matched (p = 0.0035 REVOLVER) \tabularnewline
Survey systematics stability & \checkmark{} Validated & Phase 3: FE $\Delta$$\gamma$ < 0.0005; LOSO 16/\allowbreak{}16 positive \tabularnewline
NGC/\allowbreak{}SGC asymmetry interpretation & \checkmark{} Validated & Phase 3b: asymmetry not significant (Wald p = 0.49--0.96), not geometric (mock p = 0.28--0.90), SGC sensitivity-limited (6--8$\times$ fewer voids, 44\% recovery power). Five independent diagnostics confirm footprint-dependent detectability. \tabularnewline
5$\sigma$ detection forecast (3,400 SNe) & \checkmark{} Validated & Phase 4: Monte Carlo with consistency checks passing at 10\textsuperscript{$-$13} \tabularnewline
Statistics-limited (not systematics-limited) & \checkmark{} Validated & Phase 4: baseline $\approx$ realistic; rank-1 systematic immunity demonstrated \tabularnewline
Hubble tension consistency (conditional; void depletion) & $\blacksquare$ Constrained & Void-depletion mapping yields H\textsubscript{local} consistent with SH0ES, conditional on local-underdensity depth. Mechanism independent of CMB; not yet tested with dedicated void-SN joint analysis. \tabularnewline
Photon--stress coupling ($\kappa$ = f\textsubscript{kick}/3 $\approx$ 0.00102) & $\blacktriangle$ Hypothesis & Derived from fundamental parameters (angular averaging of strain tensor). Below Planck sensitivity. Paper MTDF\_\allowbreak{}04 frames as speculative. \tabularnewline
Rotation curve phenomenology & $\blacksquare$ Constrained & Included in V18 scalars. Not independently reproduced within this programme. \tabularnewline
Redshift transition structure (z\textsubscript{cut} scan) & \checkmark{} Validated & Phase 6A: GLS z-score 3.62$\sigma$ at z\textsubscript{cut}=0.030; global scan-corrected p < 0.005. Spearman confirms (3.1$\sigma$). \tabularnewline
Weak lensing $\times$ environment (KiDS) & $\blacksquare$ Constrained & Phase 6B/\allowbreak{}B2: systematics-clean nulls with explicit 95\% upper limits. Bounded but not yet discriminating at KiDS sensitivity. \tabularnewline
Derived parameter consistency & \checkmark{} Validated & Phase 6C: all within 2.4$\sigma$ (v1-posterior-based; S\textsubscript{8} element superseded, see the row below). \tabularnewline
CMB lensing $\times$ voids & $\blacksquare$ Constrained & Phase 6D/\allowbreak{}E: low-z null (expected), BOSS benchmark recovers void-specific structure. Pipeline validated. \tabularnewline
S\textsubscript{8} cross-probe comparison & $\times$ Retired (as run) & Phase 6C2 compared the v1 Phase-5 posterior against WL surveys; that posterior is void (gauge artifact). Under the sealed v2 chains the direction reverses (enhanced growth, S8-WORSENED); quantitative revisit deferred to FWD-S8. \tabularnewline
f$\sigma$\textsubscript{8} $\times$ environment & $\blacksquare$ Constrained & Phase 6C4: $\mu$(a) prediction vs RSD + void growth. Max 0.22$\sigma$ discrimination; current data sensitivity-limited (\ {}10$\times$ improvement needed). \tabularnewline
BAO residual structure & $\blacksquare$ Constrained & Phase 6C3: $\Delta$W RMS = 0.27\%. Sound horizon shift +0.11\% absorbed into parameter degeneracies. Not detectable by any foreseeable survey. \tabularnewline
Cluster M\textsubscript{dyn}/M\textsubscript{lens} & $\blacksquare$ Constrained & Phase 6C5/\allowbreak{}C5b: $\eta$ = 1.000000 confirmed numerically from class\_\allowbreak{}mtdf transfer functions. Case B applies: M\textsubscript{dyn}/M\textsubscript{lens} = 1. Mass ratio channel closed as discriminator; growth-dependent cluster tests remain open. \tabularnewline
Weak lensing P(k) cross-correlation & $\times$ Falsifier & Not yet executed. Requires full P(k) from class\_\allowbreak{}mtdf at non-linear scales. \tabularnewline
Merging cluster offsets (Bullet Cluster) & $\times$ Falsifier & Not yet executed. Requires N-body simulations. \tabularnewline
\end{xltabular}
\endgroup

\subsection{Framework narrative}
\begin{mdframed}[linecolor=accent,linewidth=0.5pt,backgroundcolor=callout-bg,leftline=true,rightline=true,topline=true,bottomline=true,innerleftmargin=8pt,innerrightmargin=8pt,innertopmargin=6pt,innerbottommargin=6pt]
MTDF is fully CMB-compatible: Planck TTTEEE cannot distinguish MTDF from $\Lambda$CDM ($\Delta$$\chi$$^{2}$~<~1). MTDF deviates from $\Lambda$CDM at late times through environment-dependent dynamics. The Hubble tension is consistent with environmental mapping (void stress depletion), conditional on the depth of the local underdensity, rather than with early-universe modification. The SN~$\times$~void signal is the primary empirical discriminator. The finite range of the stress-field response concentrates observable effects at low redshift, consistent with the CMB being unaffected.
\end{mdframed}

\subsection{What survives intact}
\begin{enumerate}
  \item \textbf{Phase 1 strict totals}: independently reproduced to four decimal places.
  \item \textbf{SN $\times$ void environment signal}: three void finders, 16/\allowbreak{}16 LOSO, redshift decay matching a priori prediction.
  \item \textbf{Hubble tension consistency (conditional)}: entirely late-time, via void stress depletion:
    
Theoretical: \( H_{\rm local}/H_0 = 1.084 \;\Rightarrow\; 67.4 \times 1.084 = 73.1 \) (SH0ES: 73.04~$\pm$~1.04), conditional on local-underdensity depth (shallow void gives \~{}70--71)
MCMC-derived (sealed v2 chains): \( H_{0,\rm global} = 67.44 \pm 0.55 \) km/\allowbreak{}s/\allowbreak{}Mpc; applying the void-mapping factor \( S_0 = 1.084 \) gives a local-environment value \( H_{0,\rm local} \approx 73.1 \) km/\allowbreak{}s/\allowbreak{}Mpc, conditional on the depth of the local underdensity. The relation \( S_0 = \alpha\sqrt{\Omega_{\rm DE}} \) is an algebraic identity, not an independent prediction (see MTDF\_\allowbreak{}01 \S{}3.3).
  \item \textbf{Late-time phenomenology}: w(z), BAO distances, growth rate, rotation curves.
  \item \textbf{Framework stability}: MTDF recovers $\Lambda$CDM at early times, deviates at late times.
  \item \textbf{Cross-probe integrity (Phase 6, 10 tests)}: redshift transition identified (3.6$\sigma$); weak lensing and CMB lensing pipelines pass strict systematics controls; derived parameters self-consistent; f$\sigma$\textsubscript{8}, BAO residuals, and cluster mass predictions computed  -  all sensitivity-limited with current data, with quantified thresholds for future surveys.
\end{enumerate}

\subsection{Joint implications for late-time cosmological tensions}
While MTDF is not tuned to resolve any single cosmological tension, the validated framework exhibits simultaneous shifts in multiple late-time observables. In particular, the Phase~5 best-fit favours a modest reduction in $\sigma$\textsubscript{8} relative to $\Lambda$CDM best-fit values (0.795 vs 0.809), consistent with the direction preferred by weak-lensing surveys (KiDS, DES).

Independently, the environment-dependent dynamics tested in Phases~3 and 4 naturally modify local expansion observables without compromising global CMB consistency. Together, these results indicate that MTDF addresses multiple late-time phenomenological tensions within a single, unified dynamical framework, without introducing additional free parameters beyond k\textsubscript{f}.

\subsection{Paper implications}
\begin{table}[H]
\centering
\footnotesize
\begin{tabularx}{\linewidth}{l>{\RaggedRight\arraybackslash}X}
\toprule
\textbf{Paper} & \textbf{Impact} \tabularnewline
MTDF\_\allowbreak{}01 (Master) & Phase 1 strengthens several claims. MTDF\_\allowbreak{}01 updated in the v1.1.6 claims audit (see CHANGELOG). \tabularnewline
MTDF\_\allowbreak{}02 (Part I, Environmental) & Signal reproduced by an independent reimplementation and non-parametric robustness checks (Phase 3). Redshift transition near z $\approx$ 0.04, identified at 3.6$\sigma$ (Phase 6A, scan-corrected p < 0.005); the merged-sample break is not sharply localised at this value. Framework's strongest empirical evidence to date. \tabularnewline
MTDF\_\allowbreak{}04 (Part III, Photon coupling) & $\kappa$ $\approx$ f\textsubscript{kick}/3 is a theory-motivated value, consistent with but not measured by FIRAS bounds, structurally related to the early field energy fraction and below Planck sensitivity. Paper frames this as speculative, which is appropriate. \tabularnewline
MTDF\_\allowbreak{}05 (Part IV, Cosmological) & Sealed v2 full-Planck results govern (Section 6.1); Phase 2 plik-lite precursor (+0.81) and the CMB+BAO summary run (+1.25, narrower data) are pre-correction bookkeeping. Hubble tension mechanism intact, conditional on local-underdensity depth. $\Delta$$\chi$$^{2}$ vs matched control complete (see Section 6.3). \tabularnewline
MTDF\_\allowbreak{}07 (This paper) & Phase 6 adds ten cross-probe tests (6A--6E, 6C2--6C5). Title updated to Phases 1--6. All validation claims current as of February 2026. \tabularnewline
\bottomrule
\end{tabularx}
\end{table}

\section{Open questions and next steps}
\label{sec:sec8}

\subsection{NGC/\allowbreak{}SGC asymmetry diagnostics (Phase 3b, complete)}
Phase 3b (Section 5b) implemented the full asymmetry diagnostic battery and is now complete. All five tests confirm that the NGC/\allowbreak{}SGC contrast is not statistically significant, not caused by survey geometry, and attributable to SGC void-catalogue sparsity. The NGC/\allowbreak{}SGC claim status in the taxonomy (Section 7.2) has been upgraded from $\blacktriangle$~Hypothesis to \checkmark{}~Validated.

\subsection{Falsifier and discriminator tests}
\begin{table}[H]
\centering
\footnotesize
\begin{tabularx}{\linewidth}{l>{\RaggedRight\arraybackslash}X>{\RaggedRight\arraybackslash}X>{\RaggedRight\arraybackslash}X>{\RaggedRight\arraybackslash}X}
\toprule
\textbf{\#} & \textbf{Test} & \textbf{Type} & \textbf{Status} & \textbf{Timescale} \tabularnewline
1 & Full Planck plik + class\_\allowbreak{}mtdf & Discriminator & \textbf{PASSED at consistency level} (sealed v2 chains; $\Delta$$\chi$$^{2}$ vs matched control complete, see Section 6.3) & v2 and control chains sealed \tabularnewline
2 & NGC/\allowbreak{}SGC diagnostic battery & Upper limit & \textbf{COMPLETE} (sensitivity-limited, not geometric) & Complete \tabularnewline
3 & Weak lensing $\times$ environment (KiDS) & Upper limit & \textbf{COMPLETE} (Phase 6B/\allowbreak{}B2: systematics-clean null) & Complete \tabularnewline
4 & CMB lensing $\times$ voids & Upper limit & \textbf{COMPLETE} (Phase 6D/\allowbreak{}E: pipeline validated) & Complete \tabularnewline
5 & S\textsubscript{8} cross-probe comparison & Coherence & \textbf{RETIRED as run} (v1-posterior-based; direction reversed under sealed v2) & Deferred to FWD-S8 \tabularnewline
6 & f$\sigma$\textsubscript{8} $\times$ environment & Forecast & \textbf{COMPLETE} (Phase 6C4: 0.22$\sigma$, sensitivity-limited) & Complete \tabularnewline
7 & BAO residual structure & Forecast & \textbf{COMPLETE} (Phase 6C3: $\Delta$W = 0.27\%, undetectable) & Complete \tabularnewline
8 & Cluster M\textsubscript{dyn}/M\textsubscript{lens} & Resolved & \textbf{RESOLVED} (Phase 6C5/\allowbreak{}C5b: $\eta$ = 1, Case B, M\textsubscript{dyn}/M\textsubscript{lens} = 1) & Complete \tabularnewline
9 & Weak lensing P(k) (non-linear) & Discriminator & Requires full P(k) from class\_\allowbreak{}mtdf & Months \tabularnewline
10 & Merging cluster offsets & Discriminator & Requires N-body simulations & Institutional \tabularnewline
\bottomrule
\end{tabularx}
\end{table}

\begin{mdframed}[linecolor=accent,linewidth=0.5pt,backgroundcolor=callout-bg,leftline=true,rightline=true,topline=true,bottomline=true,innerleftmargin=8pt,innerrightmargin=8pt,innertopmargin=6pt,innerbottommargin=6pt]
\textbf{Phase 2 scope note, resolved by Phase 5.}
Phase 2 used a perturbative correction layer on a $\Lambda$CDM emulator, not the full modified Boltzmann equations. The full calculation now on record is the sealed v2 chain (Section 6.1), which confirms Phase 2's structural finding (CMB-level consistency; k\textsubscript{f} unconstrained by Planck) under the corrected implementation; Phase 2's and the v1 Phase 5's numerical $\Delta$$\chi$$^{2}$ values are historical.
\end{mdframed}

\section{Appendix: Technical notes}
\label{sec:sec9}

\subsection{A.1 \,k\textsubscript{f} normalisation conventions}
Three different \( k_f \) conventions were identified across the MTDF codebase during Phase 2. Understanding these is essential for interpreting historical results.

\begin{table}[H]
\centering
\footnotesize
\begin{tabularx}{\linewidth}{l>{\RaggedRight\arraybackslash}X>{\RaggedRight\arraybackslash}X>{\RaggedRight\arraybackslash}X}
\toprule
\textbf{Convention} & \textbf{"Full MTDF" k\textsubscript{f}} & \textbf{f\textsubscript{kick}} & \textbf{Source} \tabularnewline
A (old CLASS, pre-correction) & 0.10 & 0.003303 & November 2025 cobaya runs \tabularnewline
B (corrected CLASS) & 1.0 & 0.003303 & Theory prediction \tabularnewline
C (Phase 2 code, initial, corrected) & 0.001326 & 0.001326 & CosmoPower correction layer (bug) \tabularnewline
\bottomrule
\end{tabularx}
\end{table}

Convention C initially used \( f_{\rm kick} = \alpha\kappa = 0.001326 \) instead of the correct \( f_{\rm kick} = \lambda_{\rm MTDF}/24 = 0.003303 \). This factor-of-2.49 error caused a 3$\times$ overcorrection in the angular scale ratio per unit \( k_f \), leading to a spurious \( H_0 \)--\( k_f \) degeneracy in the first MCMC run (best-fit \( H_0 = 55.6 \), which is unphysical). After correction, all runs show \( H_0 \approx 67 \) and no degeneracies.

\subsection{A.2 \,V18 parameter reference}
\begin{table}[H]
\centering
\footnotesize
\begin{tabularx}{\linewidth}{l>{\RaggedRight\arraybackslash}X>{\RaggedRight\arraybackslash}X>{\RaggedRight\arraybackslash}X}
\toprule
\textbf{Parameter} & \textbf{Value} & \textbf{Uncertainty} & \textbf{Role} \tabularnewline
$\alpha$ & 1.30 & 0.26 & Stress-field coupling (void dynamics: KBC 2013 + Hoscheit 2018; counter: Kenworthy 2019) \tabularnewline
$\beta$ & 22.7 Mpc & 0.09$\times$10$^{2}$$^{3}$ m & Strain-lever length scale (void radius distribution: Pan+2012; Hamaus+2014) \tabularnewline
$\tau$ & 13.0 Gyr & 0.2 & Field age \tabularnewline
$\beta$\textsubscript{eos} & 0.573 & 0.012 & Equation-of-state parameter \tabularnewline
E & 9.1$\times$10\textsuperscript{$-$10} Pa & 0.4$\times$10\textsuperscript{$-$10} & Field energy scale \tabularnewline
$\delta$\textsubscript{bf} & 0.150 & 0.030 & Fixed post hoc from cluster total baryon fractions at M\textsubscript{500} $\sim$ 10\textsuperscript{14} M\textsubscript{$\odot$}; a first-principles derivation is not claimed (Giodini+2009; Gonzalez+2013; see MTDF\_\allowbreak{}01 \S{}5.1) \tabularnewline
z\textsubscript{t} & 0.74 & 0.10 & Transition redshift \tabularnewline
$\kappa$ & 0.00102 & 0.00003 & Photon--stress coupling (theory-motivated: $\kappa$ $\approx$ f\textsubscript{kick}/3; below FIRAS sensitivity) \tabularnewline
H\textsubscript{0} & 70.0 km/\allowbreak{}s/\allowbreak{}Mpc & 1.5 & Hubble constant (Freedman 2021 TRGB: 69.8 $\pm$ 0.6) \tabularnewline
f\textsubscript{kick} & 0.003303 & -- & Derived: (1$-$$\beta$\textsubscript{eos})$^{2}$/((1+$\alpha$)$\times$24) \tabularnewline
\bottomrule
\end{tabularx}
\end{table}

\subsection{A.3 \,Historical context: November 2025 runs}
Previous cobaya + class\_\allowbreak{}mtdf runs (November 2025) provide useful reference points. A single-point evaluation at the Planck $\Lambda$CDM best-fit gave $\chi$$^{2}$~=~614.23 (consistent with full plik TT+TE+EE). An MCMC with Planck TT-only (16 walkers, 200 steps) found \( k_f \approx 0.03 \) in the old normalisation (Convention A), with $\chi$$^{2}$ indistinguishable from $\Lambda$CDM. These early results are consistent with the current Phase 2 conclusion that Planck cannot constrain the EFE amplitude.

\section{Data and code availability}
The complete MTDF reproducibility package, including the modified CLASS solver
(\texttt{class\_\allowbreak{}mtdf}), the V19 strict validation workbook, the
\texttt{{\small\protect\path{run_validate.py}}} pipeline, the gravity-sector artefact manifest,
and the Phase 2 CosmoPower emulator cache, is archived on Zenodo
(concept DOI \href{https://doi.org/10.5281/zenodo.19741058}{{\small\protect\path{10.5281/zenodo.19741058}}};
this release, \texttt{v1.1.7}, DOI
\href{https://doi.org/10.5281/zenodo.21525855}{{\small\protect\path{10.5281/zenodo.21525855}}})
and mirrored on GitHub at
\href{https://github.com/IMesche/MTDF}{github.com/\allowbreak{}IMesche/\allowbreak{}MTDF}.
The \texttt{v1.1.7} Git tag corresponds to the archived Zenodo release. All
numerical results reported here can be regenerated from the shipped workbook
and scripts following the procedure in the top-level \texttt{{\small\protect\path{README.md}}}.

\section{References}
{\scriptsize
\begin{thebibliography}{19}
\bibitem{ref1} Aghanim, N. et al. (Planck Collaboration) 2020, "Planck 2018 results. V. Power spectra and likelihoods", \emph{A\&A}, 641, A5. DOI: \href{https://doi.org/10.1051/0004-6361/201936386}{{\small\protect\path{10.1051/0004-6361/201936386}}}
\bibitem{ref2} Aghanim, N. et al. (Planck Collaboration) 2020, "Planck 2018 results. VI. Cosmological parameters", \emph{A\&A}, 641, A6. DOI: \href{https://doi.org/10.1051/0004-6361/201833910}{{\small\protect\path{10.1051/0004-6361/201833910}}}
\bibitem{ref3} DESI Collaboration 2025, "DESI 2024 VI: Cosmological Constraints from BAO", \emph{JCAP}, 2025, 021. DOI: \href{https://doi.org/10.1088/1475-7516/2025/02/021}{{\small\protect\path{10.1088/1475-7516/2025/02/021}}}
\bibitem{ref4} Brout, D. et al. 2022, "The Pantheon+ Analysis: Cosmological Constraints", \emph{ApJ}, 938, 110. DOI: \href{https://doi.org/10.3847/1538-4357/ac8e04}{{\small\protect\path{10.3847/1538-4357/ac8e04}}}
\bibitem{ref5} Melia, F. \& Maier, R. S. 2013, "Cosmic chronometers in the R\_\allowbreak{}h=ct universe", \emph{MNRAS}, 432, 2669. DOI: \href{https://doi.org/10.1093/mnras/stt1082}{{\small\protect\path{10.1093/mnras/stt1082}}}
\bibitem{ref6} Alam, S. et al. (eBOSS Collaboration) 2021, "Completed SDSS-IV eBOSS: Cosmological implications", \emph{Phys. Rev. D}, 103, 083533. DOI: \href{https://doi.org/10.1103/PhysRevD.103.083533}{{\small\protect\path{10.1103/PhysRevD.103.083533}}}
\bibitem{ref7} Spurio Mancini, A. et al. 2022, "CosmoPower: emulating cosmological power spectra for accelerated Bayesian inference", \emph{MNRAS}, 511, 1771. DOI: \href{https://doi.org/10.1093/mnras/stac064}{{\small\protect\path{10.1093/mnras/stac064}}}
\bibitem{ref8} Torrado, J. \& Lewis, A. 2021, "Cobaya: Code for Bayesian Analysis of hierarchical physical models", \emph{JCAP}, 2021, 057. DOI: \href{https://doi.org/10.1088/1475-7516/2021/05/057}{{\small\protect\path{10.1088/1475-7516/2021/05/057}}}
\bibitem{ref9} Foreman-Mackey, D. et al. 2013, "emcee: The MCMC Hammer", \emph{PASP}, 125, 306. DOI: \href{https://doi.org/10.1086/670067}{{\small\protect\path{10.1086/670067}}}
\bibitem{ref10} Blas, D., Lesgourgues, J. \& Tram, T. 2011, "The Cosmic Linear Anisotropy Solving System (CLASS). Part II", \emph{JCAP}, 2011, 034. DOI: \href{https://doi.org/10.1088/1475-7516/2011/07/034}{{\small\protect\path{10.1088/1475-7516/2011/07/034}}}
\bibitem{ref11} Asgari, M. et al. (KiDS Collaboration) 2021, "KiDS-1000 Cosmology: Cosmic shear constraints", \emph{A\&A}, 645, A104. DOI: \href{https://doi.org/10.1051/0004-6361/202039070}{{\small\protect\path{10.1051/0004-6361/202039070}}}
\bibitem{ref12} Abbott, T. M. C. et al. (DES Collaboration) 2022, "Dark Energy Survey Year 3 results: Cosmological constraints", \emph{Phys. Rev. D}, 105, 023520. DOI: \href{https://doi.org/10.1103/PhysRevD.105.023520}{{\small\protect\path{10.1103/PhysRevD.105.023520}}}
\bibitem{ref13} Dalal, R. et al. (HSC Collaboration) 2023, "Hyper Suprime-Cam Year 3 Results: Cosmology from Cosmic Shear", \emph{Phys. Rev. D}, 108, 123519. DOI: \href{https://doi.org/10.1103/PhysRevD.108.123519}{{\small\protect\path{10.1103/PhysRevD.108.123519}}}
\bibitem{ref14} Keenan, R. C., Barger, A. J. \& Cowie, L. L. 2013, "Evidence for a \~{}300 Megaparsec Scale Under-density", \emph{ApJ}, 775, 62. DOI: \href{https://doi.org/10.1088/0004-637X/775/1/62}{{\small\protect\path{10.1088/0004-637X/775/1/62}}}
\bibitem{ref15} Pan, D. C. et al. 2012, "Cosmic Voids in SDSS DR7", \emph{MNRAS}, 421, 926. DOI: \href{https://doi.org/10.1111/j.1365-2966.2011.20197.x}{{\small\protect\path{10.1111/j.1365-2966.2011.20197.x}}}
\bibitem{ref16} Freedman, W. L. 2021, "Measurements of the Hubble Constant: Tensions in Perspective", \emph{ApJ}, 919, 16. DOI: \href{https://doi.org/10.3847/1538-4357/ac0e95}{{\small\protect\path{10.3847/1538-4357/ac0e95}}}
\bibitem{ref17} McGaugh, S. S., Lelli, F. \& Schombert, J. M. 2016, "Radial Acceleration Relation in Rotationally Supported Galaxies", \emph{PRL}, 117, 201101. DOI: \href{https://doi.org/10.1103/PhysRevLett.117.201101}{{\small\protect\path{10.1103/PhysRevLett.117.201101}}}
\bibitem{ref18} Lelli, F. et al. 2017, "One Law to Rule Them All: The Radial Acceleration Relation", \emph{ApJ}, 836, 152. DOI: \href{https://doi.org/10.3847/1538-4357/836/2/152}{{\small\protect\path{10.3847/1538-4357/836/2/152}}}
\bibitem{ref19} Riess, A. G. et al. 2022, "A Comprehensive Measurement of the Local Value of the Hubble Constant", \emph{ApJL}, 934, L7. DOI: \href{https://doi.org/10.3847/2041-8213/ac5c5b}{{\small\protect\path{10.3847/2041-8213/ac5c5b}}}
\end{thebibliography}
}


\end{document}
